\documentclass[output=paper,colorlinks,citecolor=brown]{langscibook}
\ChapterDOI{10.5281/zenodo.19708124}

\title[Speaking and acting with care]{Speaking and acting with care: The grammar of circumspection in Upper Tanana Dene}
\author{Olga Lovick\affiliation{University of Saskatchewan}}
\abstract{Upper Tanana Dene has two constructions falling into the domain of circumspection: avertive clause linkings and urgent requests. In avertive linkings, the avertive clause expresses an undesirable consequence that can be averted by the action expressed in the main clause. Somewhat unusual typologically, avertive linkings with a request as main clause are rare in Upper Tanana. I argue in this paper that this is likely due to the fact that requests are generally problematic in \ili{Northern Dene}, and that spelling out the undesirable consequences to be avoided could suggest to the addressee that their cultural knowledge is inadequate. Instead, there is a dedicated strategy for urgent requests resulting from an insubordinate conditional construction, which draws attention to the undesirable consequences of non-compliance without spelling them out explicitly. Together, these two constructions reflect the circumspect mindset of the Upper Tanana Dene, which requires them to act and speak with care.}

\IfFileExists{../localcommands.tex}{
 \addbibresource{../localbibliography.bib}
 \usepackage{tabularx,multicol}
\usepackage{url}
\urlstyle{same}

\usepackage{langsci-optional}
\usepackage{langsci-lgr}
\usepackage{langsci-gb4e}

\babelprovide[import]{hebrew}

\usepackage{paralist} %%may not be needed, currently used in questionnaire


%Siena added packages for introduction; not sure they are needed
\usepackage[normalem]{ulem}
\useunder{\uline}{\ul}{}

%from Treis:
\usepackage{listings}
\lstset{basicstyle=\ttfamily,tabsize=2,breaklines=true}

%%from Francez
% % % \usepackage{cjhebrew}
%\usepackage{tipa}  Also included by D&A, need to check with LSP whether to change it to language science tipa
\usepackage{leipzig}
%%from Dabkowski&AnderBois:

\usepackage[shortlabels]{enumitem} 
\usepackage{centernot}
\usepackage{arydshln}
\usepackage{newunicodechar}
\usepackage{csquotes}
%\let\sups\relax \usepackage{tipa} commented out on 7/3/25 Martina
% \usepackage[all]{nowidow}
\usepackage{listings} \lstset{basicstyle=\ttfamily}
\usepackage{multirow}
\usepackage{rotating}
\usepackage{pdflscape}
\usepackage{xspace}
%\usepackage{tabularx}
%\usepackage{multicol}
\usepackage{supertabular}
\usepackage{hhline}
\usepackage{bbding} %this package is for the checkmarks in Table 6 in Browne et al.

%\usepackage{qtree}
%\usepackage{tree-dvips}

\usepackage{array}
\usetikzlibrary{calc,tikzmark,matrix}
\usepgflibrary{shadings}
\usepackage{pifont}
%from Fedotov:

\usepackage{langsci-textipa}
%\usepackage{multirow}


%\usepackage{xeCJK}

%%%%%%%%%%%%%%%%%%%%%%%%%%%%%%%%%%%%%%%%%%%%%%%%%%%%
%%% EXAMPLES %%%%%%%%%%%%%%%%%%%%%%%%%%%%%%%%%%%%%%%
%%%%%%%%%%%%%%%%%%%%%%%%%%%%%%%%%%%%%%%%%%%%%%%%%%%% 

% if you want the source line of examples to be in
% italics, uncomment the following line
\renewcommand{\exfont}{\itshape}
% to add additional information to the right of
% examples, uncomment the following line

%\usepackage{langsci/styles/jambox}


%%% FOOTNOTE EXAMPLE NUMBERING
% \makeatletter
% \patchcmd{\@footnotetext}{\setcounter{fnx}{0}}{\renewcommand{\thexnumi}{\roman{xnumi}}}{}{}
% \apptocmd{\@footnotetext}{
%     \@noftnotetrue
%     \renewcommand{\thexnumi}{\arabic{xnumi}}%
% }{}{}
%
% \makeatother
\usepackage[linguistics,edges]{forest}
\usepackage{hyphenat}
\usepackage{langsci-branding}

 \makeatletter
\let\thetitle\@title
\let\theauthor\@author
\makeatother

\newcommand{\togglepaper}[1][0]{
   \bibliography{../localbibliography}
   \papernote{\scriptsize\normalfont
     \theauthor.
     \titleTemp.
    To appear in:
    Martina Faller,  Marine Vuillermet \&  Eva Schultze-Berndt (eds.). Apprehensional Constructions in a cross-linguistic perspective. Berlin: Language Science Press. [preliminary page numbering]
   }
   \pagenumbering{roman}
   \setcounter{chapter}{#1}
   \addtocounter{chapter}{-1}
}

\newbool{bookcompile}
\booltrue{bookcompile}
\newcommand{\bookorchapter}[2]{\ifbool{bookcompile}{#1}{#2}}


\newfontfamily\FedotovBoldEmptySetFont{LibertinusSerif-Italic.otf}[AutoFakeBold=2.5]
\newcommand{\FedotovBoldEmpty}{{\FedotovBoldEmptySetFont\bfseries ∅}}

%From Dabkowski&AnderBois


\let\sc\scshape
\let\rm\rmfamily
%\let\it\itshape
\let\tt\ttfamily
\let\nr\normalfont

\newcommand{\wsc}[1]{\textls[80]{\scshape#1}}

\let\textai\textit
% \let\textlc\spacedlowsmallcaps
% \let\textac\spacedallcaps
\let\textlc\textsc
\let\textac\textsc


%%%%%%%%%%%%%%%%%%%%%%%%%%%%%%%%%%%%%%%%%%%%%%%%%%%%
%%% FLAG %%%%%%%%%%%%%%%%%%%%%%%%%%%%%%%%%%%%%%%%%%%
%%%%%%%%%%%%%%%%%%%%%%%%%%%%%%%%%%%%%%%%%%%%%%%%%%%%

% \newcommand{\scott}  [1]{\textcolor{violet}{#1}}
% \newcommand{\maks}   [1]{\textcolor{blue}{#1}}
% % \newcommand{\data}   [1]{\textcolor{teal}{#1}}
% \newcommand{\new}   [1]{#1}
% \newcommand{\newnew}   [1]{\textcolor{teal}{#1}}
% \newcommand{\rev}   [1]{\textcolor{magenta}{#1}}



%%%%%%%%%%%%%%%%%%%%%%%%%%%%%%%%%%%%%%%%%%%%%%%%%%%%
%%% TABLES %%%%%%%%%%%%%%%%%%%%%%%%%%%%%%%%%%%%%%%%%
%%%%%%%%%%%%%%%%%%%%%%%%%%%%%%%%%%%%%%%%%%%%%%%%%%%%

%\newcolumntype{Y}{>{\arraybackslash}p{25.3543464286pt}} % length of L

% \newcolumntype{K}{>{\arraybackslash}p{24.25pt}}
% \newcolumntype{V}{>{\arraybackslash}p{(\textwidth/9-24pt)/2}}
%
\let\ch\relax \newcommand{\ch}[2]{\multicolumn{#1}{c}{\sc#2}}
\let\sx\relax \newcommand{\sx}[2]{\multicolumn{#1}{l}{\emph{-#2}}} % suffix
\let\su\relax \newcommand{\su}{\sx{1}} % suffix unicolumnal
\let\sd\relax \newcommand{\sd}{\sx{2}} % suffix dicolumnal
\let\cx\relax \newcommand{\cx}[2]{\multicolumn{#1}{l}{\emph{=#2}}} % clitic
\let\cu\relax \newcommand{\cu}{\cx{1}} % clitic unicolumnal
\let\cd\relax \newcommand{\cd}{\cx{2}} % suffix dicolumnal
\let\gx\relax \newcommand{\gx}[2]{\multicolumn{#1}{l}{\sc #2}} % gloss
\let\gd\relax \newcommand{\gd}{\gx{2}} % gloss dicolumnal
\let\gr\relax \newcommand{\gr}[1]{\multicolumn{2}{c}{\multirow{2}{*}{\textcolor{gray}{\sc#1}}}}
\let\db\relax \newcommand{\db}[1]{\multicolumn{2}{c}{\multirow{2}{*}{\sc#1}}}



%%%%%%%%%%%%%%%%%%%%%%%%%%%%%%%%%%%%%%%%%%%%%%%%%%%%
%%% EXAMPLES %%%%%%%%%%%%%%%%%%%%%%%%%%%%%%%%%%%%%%%
%%%%%%%%%%%%%%%%%%%%%%%%%%%%%%%%%%%%%%%%%%%%%%%%%%%%

%%% AFFIX BOUNDARY %%%%%%
% \DeclareRobustCommand{\longmn}{\textminus}
% \DeclareRobustCommand{\mn}{-}

%%% CLITIC BOUNDARY %%%%%%
% \DeclareRobustCommand{\longeq}{=}
% \makeatletter\DeclareRobustCommand{\eq}{%
%     \settowidth{\@tempdima}{-}% width of a hyphen
%     \resizebox{\@tempdima}{\height}{=}%
% }\makeatother

%%% FUZZY CLITIC %%%%%%
% \DeclareRobustCommand{\longfc}{%
%     $\overset{\text{\smash{}}}{\text{$\approx$}}$%
% }
% \makeatletter\DeclareRobustCommand{\fc}{%
%     \settowidth{\@tempdima}{-}% width of a hyphen
%     \resizebox{\@tempdima}{\height}{%
%         $\overset{\text{\smash{}}}{\text{$\approx$}}$%
%     }%
% }\makeatother

%%% REDUPLICATION %%%%%%%
% \DeclareRobustCommand{\longtl}{$\sim$}
% \makeatletter\DeclareRobustCommand{\tl}{%
%     \settowidth{\@tempdima}{-}% width of a hyphen
%     \resizebox{\@tempdima}{\height}{$\sim$}%
% }\makeatother

% \newcommand{\mn}          {\textminus}
\newcommand{\src}    [1]  {\jambox[0pt]{\rm(#1)}}
% \newcommand{\ai}     [2]  {\emph{#1} `\textsc{#2}#3'\xspace}
\newcommand{\ai}     [2]  {\emph{#1} `\textsc{#2}'\xspace}
\newcommand{\sane}   [1][]{\ai{=sa'ne}{appr}{#1}}
\newcommand{\srcn}   [1][]{\ai{kûintsû}{srcn}{#1}}
% \newcommand{\mbekane}[1][]{\ai{\eq mbe kañe}{neg.adv aux.inf}{#1}}

% \let\sc\relax \newcommand{\sc}{\normalfont\scshape}
% \let\rm\relax \newcommand{\rm}{\normalfont\rmfamily}
% \let\it\relax \newcommand{\it}{\normalfont\itfamily}

\newcommand{\lb}{{\rm[}}
\newcommand{\rb}{{\rm]}}

\newcommand{\bad}{\makebox[0pt][r]{*\,}}
\newcommand{\infel}{\makebox[0pt][r]{\#\,}}


% \newcommand{\lsqz}[1]{#1}
% \newcommand{\lsqz}[1]{\makebox[0pt][l]{#1}}

% \newcommand{\lsqu}[1]{\makebox[0pt][l]{#1}}
% \newcommand{\rsqu}[1]{\makebox[0pt][r]{#1}}
% \newcommand{\astr}{\makebox[0pt][r]{{\nr*}}}
% \newcommand{\hash}{\makebox[0pt][r]{{\nr\#}}}
% \newcommand{\qstn}{\makebox[0pt][r]{{\nr?}}}
% \newcommand{\bad}{\makebox[0pt][r]{*}}
% \newcommand{\unhappy}{\makebox[0pt][r]{\#}}


%%%%%%%%%%%%%%%%%%%%%%%%%%%%%%%%%%%%%%%%%%%%%%%%%%%%
%%% UNICODE CHARACTERS  %%%%%%%%%%%%%%%%%%%%%%%%%%%%
%%%%%%%%%%%%%%%%%%%%%%%%%%%%%%%%%%%%%%%%%%%%%%%%%%%%

\newunicodechar{⟨}{⟨}
\newunicodechar{⟩}{⟩}

\newcommand{\infix}{⟨}
\newcommand{\outfix}{⟩}



%%%%%%%%%%%%%%%%%%%%%%%%%%%%%%%%%%%%%%%%%%%%%%%%%%%%
%%% IPA %%%%%%%%%%%%%%%%%%%%%%%%%%%%%%%%%%%%%%%%%%%%
%%%%%%%%%%%%%%%%%%%%%%%%%%%%%%%%%%%%%%%%%%%%%%%%%%%%

% \newcommand{\ipa}{\textipa}

\newcommand{\sub}{\textsubscript}
% \let\super\relax \newcommand{\super}{\textsuperscript}

\newcommand{\ts}{\texttslig}
% \let\dz\relax \newcommand{\dz}{\textdzlig}
\newcommand{\tS}{\textteshlig}
\newcommand{\dZ}{\textdyoghlig}
\newcommand{\cj}{\textctc}
\newcommand{\sj}{\textctc}
\newcommand{\zj}{\textctz}
\newcommand{\tcj}{\texttctclig}
\newcommand{\tsj}{\texttctclig}
\newcommand{\dzj}{\textdctzlig}
\newcommand{\gj}{\textbardotlessj}
% \let\nj\relax \newcommand{\nj}{\textltailn}
% \newcommand{\W}{\textturnmrleg}
\newcommand{\wl}{\textturnmrleg}
% \newcommand{\1}{\ipabar{\~\i}{.5ex}{1.1}{}{}}
\newcommand{\dotlessibar}{\ipabar{\i}{.5ex}{1.1}{}{}}
\newcommand{\darkl}{\textltilde}
% \newcommand{\n}{\textsubarch}
% \newcommand{\x}{\textovercross}
\newcommand{\lowered}{\textlowering}
\newcommand{\raised}{\textraising}
\newcommand{\retracted}{\textsubbar}
\newcommand{\unreleased}{\textcorner}
% \newcommand{\syllabic}{\s}
\newcommand{\implosivet}{\texthtt}
\newcommand{\implosived}{\texthtd}
\newcommand{\implosivep}{\texthtp}
\newcommand{\implosiveb}{\texthtb}
\newcommand{\implosivek}{\texthtk}
\newcommand{\implosiveg}{\texthtg}



%%%%%%%%%%%%%%%%%%%%%%%%%%%%%%%%%%%%%%%%%%%%%%%%%%%%
%%% OTHER %%%%%%%%%%%%%%%%%%%%%%%%%%%%%%
%%%%%%%%%%%%%%%%%%%%%%%%%%%%%%%%%%%%%%%%%%%%%%%%%%%%

\newcommand{\ie}{i.\,e.\xspace}
\newcommand{\Ie}{I.\,e.\xspace}
\newcommand{\eg}{e.\,g.\xspace}
\newcommand{\Eg}{E.\,g.\xspace} 




\newcommand{\francoissource}[1]{\hfill\textnormal{{\footnotesize #1} }}
\newcommand{\E}{Ɛ\kern-1.5pt\raisebox{1pt}{̏}\kern1.5pt}
\newcommand{\ù}{{ṵ}\kern-2pt{̀}\kern2pt}
\newcommand{\ȕ}{{ṵ}\kern-2pt{̏}\kern2pt}

\newcommand{\OO}{ɔ\kern-1pt\raisebox{-1pt}{̰}\kern1pt}
\newcommand{\Ǒ}{{{ɔ̌\kern-1pt\raisebox{-.5pt}{̰}\kern1pt}}}


\newcommand{\à}{{a̰}\kern-2pt{̀}\kern2pt}
\newcommand{\ilit}[1]{#1\il{#1}}

 %% hyphenation points for line breaks
%% Normally, automatic hyphenation in LaTeX is very good
%% If a word is mis-hyphenated, add it to this file
%%
%% add information to TeX file before \begin{document} with:
%% %% hyphenation points for line breaks
%% Normally, automatic hyphenation in LaTeX is very good
%% If a word is mis-hyphenated, add it to this file
%%
%% add information to TeX file before \begin{document} with:
%% %% hyphenation points for line breaks
%% Normally, automatic hyphenation in LaTeX is very good
%% If a word is mis-hyphenated, add it to this file
%%
%% add information to TeX file before \begin{document} with:
%% \include{localhyphenation}

\hyphenation{
par-a-digm
com-ple-men-ta-tion
anaph-o-ra
Dor-drecht
Cu-shi-tic
be-ne-dic-tive
pa-ra-digms
Ethi-o-pi-an
hoch-land-ost-ku-schi-ti-schen 
Duu-raa-me
Pa-ken-dorf
Lich-ten-berk
affri-ca-te
affri-ca-tes
com-ple-ments
Eng-lish
}




\hyphenation{
par-a-digm
com-ple-men-ta-tion
anaph-o-ra
Dor-drecht
Cu-shi-tic
be-ne-dic-tive
pa-ra-digms
Ethi-o-pi-an
hoch-land-ost-ku-schi-ti-schen 
Duu-raa-me
Pa-ken-dorf
Lich-ten-berk
affri-ca-te
affri-ca-tes
com-ple-ments
Eng-lish
}




\hyphenation{
par-a-digm
com-ple-men-ta-tion
anaph-o-ra
Dor-drecht
Cu-shi-tic
be-ne-dic-tive
pa-ra-digms
Ethi-o-pi-an
hoch-land-ost-ku-schi-ti-schen 
Duu-raa-me
Pa-ken-dorf
Lich-ten-berk
affri-ca-te
affri-ca-tes
com-ple-ments
Eng-lish
}



 \boolfalse{bookcompile}
 \togglepaper[23]%%chapternumber
}{}




\begin{document}
\maketitle

\section{Introduction}
\label{section:intro}

In this paper I describe two constructions falling into the domain of circumspection in Upper Tanana, a \ili{Dene} (Athabascan) language spoken in eastern interior Alaska and the western Yukon. 
\textit{Avertive linkings} are adverbial clause linkings where the subordinate clause expresses an undesirable event that can be averted or prevented by taking the action expressed in the main clause. I choose the label \textit{avertive} over that of \textit{negative purposive}, since in Upper Tanana, this clause type differs in important ways from negated purposive clauses, as discussed in \sectref{subsection:alternative-constru}. \textit{Urgent requests} are a special type of request signaled by the presence of the marker \textit{dè'} `if, when.\textsc{fut}, \textsc{ur}'. This particle suggests that this construction arose via insubordination of a conditional clause construction. Non-compliance with these requests usually has dire consequences, but these are never spelled out explicitly. I will also show how the existence and use of these constructions reflect the circumspect mindset of the \ili{Upper Tanana Dene}.

In \sectref{subsection:careful-action}, I provide cultural background on the Upper Tanana, focusing in particular on the need for circumspect action and careful speaking. \sectref{subsection:upper-tanana} provides linguistic background on the \ili{Upper Tanana} language. The avertive clause linking is described in detail in \sectref{section:avertive}, urgent requests in \sectref{section:apprehensional}. \sectref{section:speaking-carefully} situates both constructions in the cultural context. The findings are summarized in \sectref{section:conclusion}.

In my discussion of cultural traits of the Upper Tanana throughout this paper, I draw on ethnographic resources on other \ili{Northern Dene} groups, in particular on the work by \citet{NelsonR1973} on \ili{Gwich'in}, \citet{NelsonRMautnerKBaneG1982} and \citet{NelsonR1983} on \ili{Koyukon}, and \citet{ScollonRScollonS1981} on \ili{Tanacross}  (all these groups are located in Alaska) as well as \citet{RushforthS1985} and \citet{RushforthSChisholmJ1991} on \ili{Bearlake} (Northwest Territories). The ethnographic literature on Upper Tanana itself focuses on topics that do not relate directly to the focus of the present paper. \citet{McKennanR1959} and \citet{McKennanR1981} provide broad cultural overviews while \citet{McKennanR1969} is concerned with population development during the contact period. The early work by \citet{GuedonMF1974, GuedonMF1981} describes the potlatch, one of the most important cultural institutions; her most recent work \citep{GuedonMF2005} focuses on the spiritual life and particularly on shamanism. 

Even though all of the above cultures differ from each other in many details, the fundamental patterns are very similar, and, as has been argued by \citet{McKennanR1981} for the groups known today as Lower, Middle, and Upper Tanana as well as \ili{Tanacross}, and by \citet{KraussM2000Koydialects} for Alaskan Dene languages, the group distinctions made by ethnographers and linguists do not reflect the perspective of the \ili{Dene} themselves, who view their communities as part of an intricate network covering a large portion of North America. In taking this approach, I follow the practice of \citet{ScollonRScollonS1981} in their work on interethnic communication, who focus on two groups they label \textit{Athabaskans} and \textit{speakers of English}. By \textit{Athabaskan}, they ``mean anyone who has been socialized to a set of communicative patterns which have their roots in the Athabaskan languages'' (p.\ 12).  By the same token, I use the term \textit{Dene} to refer to anyone socialized into any (Northern) Dene culture, while also paying attention to customs specific to the Upper Tanana area. 

It is also important to note that many of the cultural traits described here are less salient to younger individuals, likely due to the influence of mainstream American culture.

\subsection{The need for circumspection: Careful speaking, careful action}
\label{subsection:careful-action}

The interactions of Alaskan Dene are characterized by circumspect behaviour, linguistic and non-linguistic. Some aspects of linguistic circumspection have been described by \citet{ScollonRScollonS1981}. Of particular importance for our purposes are two points raised there: reluctance to talk about the future (coupled with a need for preparedness) and the principle of non-interference.

For an Alaskan Dene person, to talk with any degree of certainty about the future constitutes a violation of \textit{įįjih}, the moral system governing an individual's behavior (cf.\  \citealt[chap.\ 3]{GuedonMF2005}; \citealt{LovickO2016, LovickOTuttleS2019} for more discussion of \ili{Upper Tanana} \textit{įįjih}; similar concepts are described for many Alaskan Dene groups). Making plans for the future, or expressing an expectation about what the future will be like (or that the future will be good), is to court bad luck, and is thus consistently avoided \citep[20]{ScollonRScollonS1981}.  \citet[194--195]{NelsonR1973} notes the unreliability of the weather in interior Alaska, and the difficulty even experienced hunters encounter when trying to predict the weather. The natural world is thought to be too large to be a safe discussion topic; \citet[45]{NelsonR1983} reports that \ili{Koyukon} children are often reminded of the need for humility in the face of nature using the admonition ``Don't talk; your mouth is too small.''

As a consequence, the \ili{Upper Tanana Dene} place a lot of cultural value on the concept of \textit{ts'edòònih}, an uninflectable word of uncertain lexical category often translated as `you don't know what's gonna happen' (author's fieldnotes) or `you never know' \citep{KariJ1997stem}.  Implicit in the concept of \textit{ts'edòònih} is the recognition that there are many things beyond an individual's control, and the awareness that it is each person's responsibility to be prepared for all of them at all times. Part of this preparedness is the need to look out for one another, as articulated by Mrs. Avis Sam (N) in a recent text about \textit{ts'edòònih} (\ref{tsedoonih}).\footnote{All data is presented in the practical orthography. Tone is not usually indicated in the practical orthography but marked in the bulk of the data in this paper, with the exception of the published data from the Scottie Creek dialect \citep{TyoneM1996}, which is presented as in the original source. High tone, associated with standard negation and several lexical items, is marked for all data from the Beaver Creek, Northway, and Tetlin dialects. Low tone originating from Proto-\ili{Dene} constriction \citep{KraussM2005} is marked for the Northway and Beaver Creek data but not for the Tetlin data, as low tone is not phonemic in that variety \citep{MinouraN1994, LovickO2020}.}

\ea 
\label{tsedoonih}
\gll Nił-'èh na-ts'i-dek dè' ts'iikeey iin ts'edòònih xà' hǫǫsú' {nił-k'à-nal-tà'. \textnormal{(N)}} \\
\textsc{recp.po}-with \textsc{prmb-1pl.s:ø.ipfv:ø-pl}.go:\textsc{ipfv:cust} when children \textsc{pl} never.know because well \textsc{recp.po}-care-\textsc{qual:ø.ipfv:2pl.s:l}-take.care:\textsc{ipfv} \\
\glt `When we walk around together, children, take good care of each other because of \textit{ts'edòònih}.' \hfill [UTOLAF19Jan0501:019]\\
\z

In addition to this general mindset of circumspection, Alaskan Dene traditionally perform numerous small, often ritualistic acts that serve to prevent undesirable events or to ensure favorable outcomes. This is illustrated in (\ref{fish-heads}), which describes the proper placement of whitefish on a drying rack in the Upper Tanana area.\footnote{Whitefish are caught during their annual migration to their spawning grounds, hence there is a `forward' and a `backward' relative to the direction of their migration.} 

\begin{exe}
\ex 
\label{fish-heads}
\begin{xlista}
\ex 
\gll A-dógn da-ts'e-łeeg tah chih k'at'eey né' nts'ą' k'a da-tsaa-dláál u-tthi' ts'ąy a-nóó ts'ą' shyį́į́' da-tsaa-łaał {u-tthi'. \textnormal{(T)}} \\
\textsc{ntrl}-up:\textsc{ar} up-\textsc{1pl.s:ø.ipfv:h}-handle.\textsc{pl.o:ipfv:cust:nmz} when also \textsc{neg} backward:\textsc{all} \textsc{advzr} \textsc{neg} up-\textsc{1pl.s:inc:aa.ipfv:ø}-keep.\textsc{pl.o:fut:neg} \textsc{3pl.psr}-head side \textsc{ntrl}-ahead:\textsc{all} \textsc{advzr} only up-\textsc{1pl.s:inc:aa.ipfv:ø}-keep.\textsc{pl.o:fut} \textsc{3pl.psr}-head \\
\glt `Also when we put [whitefish] up [on the drying rack], we may not keep it up there backward, we have to keep it there with the sides of the head facing forward, their heads.' \hfill [UTOLVDN11Jul2903:003]\\
\ex 
\gll Ay tah chih k'at'eey keł na-tat-dlą́ą́ he-ndiik na-né' nts'ą' da-hiiy-eh-łeek tah. \\
\textsc{3sg} when also \textsc{neg} ? \textsc{it-3pl.s:inc:aa.ipfv:d}-be.many:\textsc{fut:neg} \textsc{3pl.s:ø.ipfv:ø}-say:\textsc{ipfv} \textsc{med}-backward:\textsc{all} \textsc{advzr} up-\textsc{3pl>3pl-ø.ipfv:ø}-handle.\textsc{pl.o:ipfv:cust} when \\
\glt `There won't be as many, they say, when they put them up facing backward.' \hfill [UTOLVDN11Jul2903:004]\\
\end{xlista}
\end{exe}

 There are countless acts such as these; \citet[chs. 3, 4, 9, 14]{NelsonRMautnerKBaneG1982} list a large number of them for the Koyukon people, but are careful to note that there are many more. Lists of these acts, and of the consequences they are intended to avert, are a rich source for avertive linkings (see also \sectref{section:speaking-carefully}).

Knowledge of these acts is part and parcel of being an \ili{Alaskan Dene}; in the Upper Tanana area, this is one component of ``knowing one's \textit{įįjih}'' (author's fieldnotes), i.e.\  having a solid grasp of cultural knowledge. A competent individual is expected to know (and ideally adhere to) their \textit{įįjih}, and \ili{Upper Tanana} speakers are very careful not to suggest that an individual's knowledge of \textit{įįjih} is deficient (cf.\  \citealt[273]{LovickO2016}).

This last point is linked to another fundamental tenet of \ili{Northern Dene} culture: non-interference with another's autonomy (cf.\  \citealp{ScollonRScollonS1981} for \ili{Alaskan Dene}, \citealp{RushforthSChisholmJ1991} for \ili{Bearlake Dene}, or \citealp{BassoK1979, BassoK1990} for \ili{Western Apache}). A direct consequence of this attitude is a general reluctance to utter questions and requests \citep{ScollonRScollonS1981}. These authors note that both types of speech act are routinely avoided, while \citet{RushforthS1985} and \citet{RushforthSChisholmJ1991} observe a marked preference for indirect requesting strategies.\footnote{This is also the reason for my terminological choice to talk about requests rather than commands.} \citet{LovickO2016} finds that, while direct affirmative requests are common in Upper Tanana, their form depends on the relationship between the interlocutors as well as on the imposition resulting from complying with the request. Negative requests are strongly avoided in \ili{Upper Tanana} (cf.\  \citealt{LovickO2016, LovickOTuttleS2019}), not merely because of the principle of non-interference, but also because they could be interpreted as suggesting that the addressee's knowledge of \textit{įįjih} is deficient.

Thus, traditional \ili{Upper Tanana} speakers are very circumspect both in words and actions. Verbally, they avoid interfering in another's autonomy, making predictions about the future, or talking about things for which their ``mouth is too small''. Practically, they value careful, often symbolic, interactions with animate and inanimate entities, as well as preparedness for all matters outside their control. 

\subsection{Upper Tanana language}
\label{subsection:upper-tanana}

\ili{Upper Tanana} is a \ili{Northern Dene} language spoken by 30–50 mostly elderly speakers in eastern interior Alaska and the Yukon Territory. \citet{MinouraN1994} identifies five dialects; this paper contains data from my own fieldwork on the Northway (N), Tetlin (T), and Beaver Creek (B) dialects as well as from a story collection \citep{TyoneM1996} in the Scottie Creek (S) dialect.

Morphologically, \ili{Upper Tanana} can be classified as polysynthetic, incorporating, and fusional. It is characterized by particularly rich verbal morphology which can be represented by the template indicating morpheme order in \tabref{tab:UT-template-general1}.

\begin{table}[htp]
\begin{tabular}{@{}lllllllllllllll@{}}
\toprule
 12  & 11 & 10 & 9    & 8      & 7    & \#& 6    & 5 & 4 & 3    & 2    &  1   & 0    & -1  \\ \midrule
\rotatebox{90}{Postpositional object} & \rotatebox{90}{Postposition} & \rotatebox{90}{Adverbial-derivational} & \rotatebox{90}{Iterative} & \rotatebox{90}{Incorporate} & \rotatebox{90}{Distributive} & \rotatebox{90}{-- Disjunct boundary --}& \rotatebox{90}{Pronominal}& \rotatebox{90}{Qualifier} & \rotatebox{90}{Conjugation} & \rotatebox{90}{Mode} &  \rotatebox{90}{Subject} & \rotatebox{90}{Voice/valence} & \rotatebox{90}{\textbf{Stem}} & \rotatebox{90}{Suffix}\\ \bottomrule
\end{tabular}
\caption{The Upper Tanana verb, represented as a template}
\label{tab:UT-template-general1}
\end{table}

\noindent The basic lexical unit in the verbal domain is the \textit{verb theme} which consists minimally of a verb stem and a voice/valence marker. Many verb themes contain additional lexical prefixes and form the basis for further derivations. There are no nonfinite verb forms; each verb is obligatorily inflected for subject (positions 2 or 6 in the template), conjugation (position 4) and mode (position 3). Transitive verb themes are inflected for direct object (position 6). Morphemes in positions 6–1 often fuse into a single syllable, which I gloss as a portmanteau here and in other work, with zero expression indicated in the gloss (see \citealt[6--10]{LovickO2020} for background on \ili{Upper Tanana} glossing specifically and \citealt{HargusSLovickOTuttleS2020} for more information on glossing conventions in \ili{Dene} languages). Other lexical classes have much more limited morphology; nouns can be inflected for possessor and (free) postpositions for postpositional object.

Upper Tanana distinguishes four \textit{modes}: imperfective, perfective, optative, and future.\footnote{The label \textit{modes} is somewhat misleading, but for lack of a better term, I continue to use it here.} The first two encode viewpoint aspect, while the latter two have modal functions. The optative expresses primarily deontic modality such as desire or obligation, while the future has mostly epistemic functions \citep{LovickO2017}. The formation of all modes involves prefixes (conjugation and mode in \tabref{tab:UT-template-general1}) as well as stem changes originating in historic suffixation \citep{LeerJ1979}; the future additionally requires the presence of the inceptive prefix (one of the qualifiers in \tabref{tab:UT-template-general1}). \ili{Upper Tanana} also has productive negative inflection as illustrated in (\ref{pos-neg}). Negative inflection is used only in standard negation and involves changes in the conjugation/mode domain as well as a negative suffix which is absorbed into the stem syllable with sometimes idiosyncratic effects; in (\ref{pos-neg}b), it voices the stem-final consonant. Standard negation also involves tonal changes; the unmarked pre-stem syllable in (\ref{pos-neg}a) becomes low in the negative form, while all stems (unmarked and low-marked alike) become high. Standard negation additionally requires the presence of the negator \textit{k'à(t'eey)} under most circumstances.

\begin{exe}
\ex \label{pos-neg}
\begin{xlista}
\ex
\gll shyì' {įį-'ààł \textnormal{(N)}}\\
meat \textsc{3sg.s:aa.pfv:ø}-eat:\textsc{pfv} \\
\glt `s/he ate the meat'  \hfill [UTCBAF13Nov0802:008]
\ex 
\gll k'àt'eey shyì' {ì-'áál \textnormal{(N)}}\\
\textsc{neg} meat \textsc{3sg.s:neg.pfv:ø}-eat:\textsc{pfv:neg} \\
\glt `s/he didn't eat the meat'  \hfill [UTCBAF13Nov0802:009] 
\end{xlista}
\end{exe}


Upper Tanana is generally verb-final, although some syntactic material may follow the verb. Arguments of verbs and postpositions as well as nominal possessors may be expressed by free noun phrases or marked on the verb, postposition, or noun.% (\ref{ut-morphosyntax}).

%\begin{exe}
%\ex \label{ut-morphosyntax}
%\begin{xlista}
%\ex
%\gll Haniign gaay xah nii-ts'et-dek tl'aan hu-k'ii-nts-eh-tsiił. (N) \\
%creek small for \textsc{term:it-1pl.s:ø.ipv:d-pl}.go:\textsc{ipv:cust} and \textsc{3pl.p}-wash-\textsc{1pl>indf:inc-ø.ipv:h}-wash:\textsc{ipv:cust} \\
%\glt `We would go to the slough and wash them.' \\
%\ex
%\gll Ts'įį-tsuul dą' {Dziin Choh Dziin'} hǫǫ-łįį tah sh-nąą nee-xah {ts'iiniin shyign'} %naa-y-ne-kįį nts'ą'. (T)\\
%\textsc{1pl.s:ø.pfv:ø}-be.small:\textsc{ipv} when:\textsc{pst} {Christmas Day} \textsc{ar.s:ø.pfv:ø}-be:\textsc{ipv} when \textsc{1sg.psr}-mother \textsc{1pl.p}-for doll \textsc{cont:it-3sg>3sg:ø.ipv:h}-sew:\textsc{ipv:cust} and \\
%\glt `When we were small, when it was Christmas Day, my mother would sew dolls for us.' \\
%\end{xlista}
%\end{exe}

Important for the discussion in this paper is the structure of adverbial clauses, which are characterized by the presence of a clause-final subordinator. Many subordinators are formally identical to postpositions and require nominalization of the verb (marked formally by a historical vocalic suffix which is absorbed into the verb stem; common effects of this suffix include voicing of stem-final consonants, deletion of stem-final \textit{h}, and lengthening of the stem vowel);\footnote{In his description of \ili{Dëne Sułiné}, \citet[375–381]{CookED2004} argues in favor of treating such clauses as nominal objects of postpositions.} many also require the verb in either clause to be in a particular mode. The order of adverbial clause and main clause is generally free but often follows iconic principles. Two types of adverbial clauses are illustrated in (\ref{adverbial}).

\begin{exe}
\ex\label{adverbial}
\begin{xlista}
\ex
\gll {Ay èh} dhįį-t'eh dè' shit {taa-'aał. \textnormal{(N)}} \\
and \textsc{dh.pfv:2sg.s:ø}-get.burned:\textsc{pfv} if scar \textsc{3sg.s:inc:aa.ipfv:ø}-classify.\textsc{co:fut} \\
\glt `And if you get burned, there will be a scar.' \\ \hfill [UTOLVDN14Apr2601:051]\\
\ex
\gll  {\ldots} ishyiit nii-ts'in-deeł na-nts'-uu-'aal {xah {\ldots} \textnormal{(N)}} \\
{} there \textsc{term:it-1pl.s:n.pfv:d-pl}.go:\textsc{pfv} \textsc{cont-1pl>indf-opt:ø}-eat:\textsc{opt:nmz} \textsc{purp} \\
\glt `{\ldots} we went back there in order to eat \ldots' \hfill [UTOLVDN14Apr2602:115] \\
\end{xlista}
\end{exe}

Example \noindent (\ref{adverbial}a) illustrates a conditional linking formed using the subordinator \textit{dè'} `if, when.\textsc{fut}', which is formally identical to the postposition \textsc{po}+\textit{dè'} `at a time \textsc{po} in the future'. This subordinator does not generally cause nominalization. The verb in the condition may be in any mode, while the verb in the consequent is almost always in the future. The condition usually, but not necessarily, precedes the consequent.  

(\ref{adverbial}b) illustrates a purposive clause formed using the subordinator \textit{xah} `in order to, \textsc{purp}', related to the postposition \textsc{po}+\textit{xah} `for \textsc{po}, because of \textsc{po}'.\footnote{This subordinator is also used in causal clauses with the meaning `because', which have a slightly different structure from the purpose clauses described here.} This subordinator usually triggers nominalization of the subordinate verb. Purpose clauses are always in the optative; there are no mode restrictions on the verb in the main clause. The main clause usually precedes the purpose clause.

Also important for this paper is the formation of (affirmative) requests, originally described in \citet{LovickO2016}. The two most common types described there are requests in the imperfective and optative modes (\ref{ipfv-opt}). 

\begin{exe}
\ex \label{ipfv-opt}
\begin{xlista}
\ex 
\gll ``Jah shyi' ay iin tl'a-'ah-łeey,'' nee-'eh-niik. \\
this meat 3 \textsc{pl} to-\textsc{ø.ipfv:2pl.s:ø}-handle:\textsc{pl.o:ipfv:imp} \textsc{1pl.o-3sg.s:ø.ipfv:h}-say:\textsc{ipfv:cust} \\
\glt `{``}Give (\textsc{pl}) them the meat here,'' he would say to us.' \\ \hfill [UTCBAF13Nov1203:039]\\
\ex 
\gll Mb-aa n-che' ski-dǫǫ-teey \ldots\\
\textsc{3sg.p}-for \textsc{2sg.psr}-tail across-\textsc{qual:opt:2sg.s:ø}-handle.\textsc{lro:opt} \\
\glt `Put your (\textsc{sg}) tail across for him \ldots' \hfill [UTOLVDN14Nov2301:359]\\
\end{xlista}
\end{exe}

Generally, requests in the imperfective are used when the speaker is comfortable making a request, either because of the relationship between the interlocutors (e.g.\  a father addressing his children as in (\ref{ipfv-opt}a), or an interaction between close friends or siblings) or because of the trivial nature of the action to be performed \citep[280--283]{LovickO2016}. Requests in the optative are used, by contrast, in situations where the request is an imposition, either because of the nature of the action to be performed or because of a less easy relationship between the speaker and the addressee. In (\ref{ipfv-opt}b), the speaker is addressing a Fox—a stranger to her. Optative requests are used in situations where the speaker does not feel entitled to making a request at all \citep[283--285]{LovickO2016}.





\subsection{Data}
\label{subsection:data}

The bulk of the data for this study comes from a corpus comprising over 8,000 utterances of mostly narrative text recorded and processed by the author between 2006--2019. About 15 speakers are represented there, but the majority of text comes from two speakers of the Tetlin dialect and four speakers of the Northway dialect. Another 2,000 or so utterances come from the story collection by Mary \citet{TyoneM1996} in the Scottie Creek dialect. I mined this corpus for instances of avertive clauses using a string search for the avertive subordinator \textit{ch'à'} and the urgent request marker \textit{dè'}. I also searched for common translations of both construction, e.g.\  ``so that {\ldots} not'',  ``to prevent that'', or ``make sure''. Selected examples were discussed with speakers during the transcription and translation process; the notes I took during those discussions also fed into this paper.

The narrative corpus is supplemented by targeted elicitation. Elicitation of avertive clauses followed the methodology outlined in \citet{VuillermetM2018questionnaire}. Using the contact language (English), I developed a scenario, e.g., walking about in the bush. I then asked for a translation of a target sentence into Upper Tanana, e.g., `I am going to boil the water so I don't get sick from it'  and discussed the result with the speaker(s): is this the only way of saying it? What other ways of expressing this are there? Why did you phrase it this way? I found that the best results were obtained when speakers could work in pairs, since they could discuss the translation with each other, but this was not always feasible. 

Direct elicitation of apprehensional requests is more challenging (see \citealt{LovickO2016, LovickOTuttleS2019} for discussion of the challenges of request elicitation), so instead, I selected examples from the corpus and discussed those with speakers, paying attention to how changing the form of the request would impact its meaning and use. 

All data collected by the author or her team is cited by its corpus identifier.

\section{Avertive linkings}
\label{section:avertive}

Avertive linkings are a type of adverbial linking. The structural properties of avertive linkings and some structural restrictions are described in \sectref{subsection:avertive-structure}. Their semantics is discussed in \sectref{subsection:avertive-semantics}; alternative constructions with similar semantics are described in \sectref{subsection:alternative-constru}. In \sectref{subsection:avert-beyond} I show that avertive linkings with similar structural properties are attested in several other \ili{Dene} languages. The origin of the avertive subordinator used in all of these languages is discussed in \sectref{subsection:avertive-origin}. This section closes with a discussion of the relative frequency of avertive linkings in different speech acts.


\subsection{Structural properties}
\label{subsection:avertive-structure}

Avertive linkings always involve the subordinator \textit{ch'à'} `\textsc{avert}, lest' in clause-final position. The verb in the avertive clause is usually in the future (\ref{undesirablea})--(\ref{undesirablec}) or in the inceptive perfective mode, which often has the meaning of `immediate future' ((\ref{undesirabled}), cf.\  \citealt{LovickO2017}). The avertive clause (bolded throughout this paper) usually follows the main clause (\ref{undesirablea})--(\ref{undesirablec}), but the opposite order is also attested~(\ref{undesirabled}). 


\begin{exe}
\ex \label{undesirable}
\begin{xlista}
\ex\label{undesirablea}
\gll Huu-ła' tah chih tl'uuł net-chüh \textbf{huu-łaa-güü} \textbf{taa-łeel} \textbf{ch'a'}. \textnormal{(S}) \\
\textsc{3pl.psr}-hand among also rope \textsc{3sg.s:qual:dh.pfv:d}-be.tied:\textsc{pfv} \textsc{3pl.psr}-finger-forked \textsc{3sg.s:inc:aa.ipfv:ø}-be:\textsc{fut:nmz} \textsc{avert} \\
\glt `Rope was tied between their fingers \textbf{to prevent them from [having] spaces between their fingers}.' \hfill \citep[19]{TyoneM1996}\\
\ex\label{undesirableb}
\gll Di-nįį' t'eey hiiy-ee-k'ia dégn' nts'ą' \textbf{na-h-nal-xon} \textbf{ch'a'} {\ldots} \textnormal{(T)} \\
\textsc{refl.psr}-face \textsc{advzr} \textsc{3pl>3pl-dh.pfv:ø}-exercise:\textsc{pfv} up:\textsc{all} \textsc{advzr} down-\textsc{3pl.s-inc:qual:aa.ipfv:l}-sag:\textsc{fut:nmz} \textsc{avert} \\
\glt `They train their faces upward (i.e.\  they wash their faces with upward motions) \textbf{to prevent them from sagging down} \ldots' \\ \hfill [UTOLVDN11Jul2901:033] \\
\ex\label{undesirablec}
\gll Ay xah \textbf{doo} \textbf{iin} \textbf{nįį-deel} \textbf{iin} \textbf{i-tįl-nüü} \textbf{ch'à'} ay xah ch'à' hu'èh {dhįįdah. \textnormal{(N)}} \\
\textsc{3sg} because who \textsc{pl} \textsc{3pl.s:n.pfv:ø-pl}.go:\textsc{pfv:nmz} \textsc{pl} \textsc{pp-inc:aa.ipfv:2sg.s:l}-forget:\textsc{fut:nmz} \textsc{avert} \textsc{3sg} because \textsc{foc} \textsc{3pl.po}-with \textsc{dh.pfv:2sg.s:ø-sg}.sit:\textsc{ipfv} \\
\glt `That's why, \textbf{so you don't forget those who came}, that is why you sit with them.' \hfill [UTOLVDN14May0508:038] 
\ex\label{undesirabled}
\gll {Ay tl'ààn} \textbf{tuu} \textbf{tèt-shyoo} \textbf{ch'à'} hà-nì-t'ah, dits'än. \textnormal{(N)}\\
and.then water \textsc{3sg.s:inc:dh.pfv:d}-get.splashed:\textsc{pfv:nmz} \textsc{avert} up/out-\textsc{3sg.s:n.ipfv:ø-sg}.fly:\textsc{ipfv} duck\\
\glt `And \textbf{to avoid getting splashed}, the duck flies up.' \\ \hfill [UTOLVDN07Aug2205:048]
\end{xlista}
\end{exe}

The main clause in an avertive linking may have any person subject: third person in (\ref{undesirablea}), (\ref{undesirableb}), (\ref{undesirabled}), second person in (\ref{undesirablec}), and first person in (\ref{first}). 

\ea \label{first}
\gll Nts'ą̀ą̀' \textbf{tat-xòl} \textbf{ch'à'} nts'ą̀ą̀' łaan {ts'uh-'įį? \textnormal{(N)}} \\
how \textsc{3pl.s:inc:aa.ipfv:d}-break:\textsc{fut:nmz} \textsc{avert} how truly \textsc{1pl.s:opt:h}-act:\textsc{opt} \\
\glt `What should we do \textbf{so [the eggs] won't break?}' \\ \hfill [UTOLVDN07Aug2205:054]
\z

Similarly, there are no person restrictions on the avertive clause's subject. While third person seems the most common ((\ref{undesirablea})--(\ref{undesirablec}), (\ref{first})), second person (\ref{undesirabled}) and first person (\ref{pov1}) are attested. There also is no co-reference requirement between the two clauses' subjects ((\ref{undesirable})--(\ref{pov1})).

\ea \label{pov1}
\gll ``Sh-taay, tl'uuł  sh-tl'ädn hi-'įh-shüh \textbf{sta-tih-haał} \textbf{ch'a'},''  {\ldots} \textnormal{(S})\\
\textsc{1sg.psr}-paternal.uncle rope \textsc{1sg.psr}-waist \textsc{pp-ø.ipfv:2sg.s:h}-tie:\textsc{ipfv} away-\textsc{inc:aa.ipfv:1sg.s:ø-sg}.go:\textsc{fut:nmz} \textsc{avert} \\
\glt {`}``My uncle, tie a rope around my waist \textbf{lest I run away},'' \ldots'\\ \hfill \citep[26]{TyoneM1996}
\z

While the avertive clause is always in the future or the inceptive perfective, there appear to be no mode restrictions on the main clause. 


\subsubsection{Structural restrictions}
\label{subsubsection:avertive-struct-restr}

The only structural restriction on avertive linkings is that the main clause has to be affirmative. Avertive linkings with negated main clauses are not attested in the corpus, nor could they be elicited.\footnote{This is true for standard negation, negative requests, and negation of existence.} Speakers consistently employ different constructions, e.g.\  (\ref{negmaina}), where the undesirable consequence of tooth decay is expressed by a causal clause, or (\ref{negmainb}), where Raven's feared theft of the sun is expressed by simple juxtaposition.

\begin{exe}
\ex \label{negmain}
\begin{xlista}
\ex\label{negmaina}
\gll  {\ldots} k'à hii-nel-déél, \textbf{hu-ts'änh} \textbf{hu-xù'} \textbf{tah-jiit} {\textbf{xah}. \textnormal{(B)}}\\
{} \textsc{neg} \textsc{3pl.s>3pl.o-qual:ø.ipfv:l}-eat:\textsc{pl.o:ipfv:neg} \textsc{3pl.po}-from \textsc{3pl.psr}-teeth \textsc{3pl:inc:aa.ipfv:h}-rot:\textsc{fut} because \\
\glt `{\ldots} they do not eat [cranberries], because their teeth will rot from them.' \hfill [UTOLVDN10Jul2602:063]
\ex\label{negmainb}
\gll Sǫ̀' ay chih èh dul-xoo, \textbf{ä́n} \textbf{y-èh} {\textbf{tthi-taa-haał}! \textnormal{(N)}} \\
\textsc{proh} \textsc{3sg} too with \textsc{3sg.s:opt:l}-play:\textsc{opt} outside:\textsc{all} \textsc{3sg.s>3sg.po}-with out-\textsc{inc:aa.ipfv:ø-sg}.go:\textsc{fut}\\
\glt `Don't let him play with [the sun], he will go outside with it!' \\ \hfill [UTOLAF18Jun1001:009]
\end{xlista}
\end{exe}

The reason for this is unclear. The small overall number of avertive linkings in the corpus (see \sectref{subsubsection:avertive-speech}) raises the possibility that this is an accidental gap, but the fact that elicitation of avertive linkings with negative or prohibitive main clauses consistently yielded constructions similar to those in (\ref{negmain}) suggests that this is a true structural restriction.


\subsection{Semantic properties}
\label{subsection:avertive-semantics}

Semantically, the avertive clause expresses an undesirable event that can be prevented by taking the action expressed in the main clause. Pubescent girls traditionally were not supposed to spread their fingers; to avoid this, their fingers were tied together (\ref{undesirablea}).\footnote{One anonymous reviewer suggests that this may be to avoid the appearance of witchcraft, which in some Dene groups involves hand gestures. To my knowledge, nothing of the sort is reported for the Upper Tanana (or related Alaskan Dene) area, but this may be due to a general reluctance to discuss such topics.} (\ref{undesirableb}) describes a measure to prevent skin aging. (\ref{undesirablec}) is about courtesy to potlatch visitors; the person putting on a potlatch has the responsibility of spending a little time with each visitor so they will remember who came (and give gifts appropriately). (\ref{undesirabled}) describes duck behavior that is exploited for egg gathering.

Undesirability is a semantic component of this construction, not a pragmatic implicature. Avertive clauses cannot be used in situations where the possible consequence is not undesirable. In most cases, the event in the avertive clause is intrinsically undesirable, but this is not required. (\ref{intrinsic}) comes from elicitation. The scenario given to the speakers was that several children are playing hide-and-seek. Two of them are sharing a hiding spot. When the seeker comes close, the girl is about to make a noise: either a scream (intrinsically undesirable; (\ref{intrinsica})) or a giggle (situationally undesirable; (\ref{intrinsicb})). The speakers originally suggested the negated purposive clause in (\ref{intrinsicc}) to express the scenario involving laughter, but decided after some discussion that the avertive clause in (\ref{intrinsicb}) was preferable. (The problematic status of negated purposive clauses in \ili{Upper Tanana} is discussed in \sectref{subsection:alternative-constru}.)\footnote{In (\ref{intrinsicb}), the speaker did not repeat the main clause, hence it is included here in brackets.}

\begin{exe}
\ex \label{intrinsic}
\begin{xlista}
\ex \label{intrinsica}
\gll D-įįłà' u-thaat da-dį̀h-'ąą \textbf{taa-säl} {\textbf{ch'à'} . \textnormal{(N)}} \\
\textsc{refl.psr}-hand \textsc{3sg.psr}-mouth closed-\textsc{3sg.s:qual:aa.pfv:h}-handle.\textsc{co:pfv} \textsc{3sg.s:inc:aa.ipfv:ø}-scream:\textsc{fut:nmz} \textsc{avert} \\
\glt `He put his hand over her mouth \textbf{so she wouldn't scream}.' \\ \hfill [UTOLAF18Jun0805:002]
\ex \label{intrinsicb}
\gll {\ob}D-įįłà' u-thaat da-dį̀h-'ąą{\cb} \textbf{taa-dlòò} {\textbf{ch'à'}. \textnormal{(N)}}\\
\textsc{refl.psr}-hand \textsc{3sg.psr}-mouth closed-\textsc{3sg.s:qual:aa.pfv:h}-handle.\textsc{co:pfv} \textsc{3sg.s:inc:aa.ipfv:ø}-laugh:\textsc{fut:nmz} \textsc{avert} \\
\glt `[He put his hand over her mouth] \textbf{so she wouldn't laugh}.' \\ \hfill [UTOLAF18Jun0805:004]
\ex \label{intrinsicc}
\gll ? D-įįłà' u-thaat da-dį̀h-'ąą \textbf{k'à} \textbf{uu-dlóó} {\textbf{xah}. \textnormal{(N)}}\\
{} \textsc{refl.psr}-hand \textsc{3sg.psr}-mouth closed-\textsc{3sg.s:qual:aa.pfv:h}-handle.\textsc{co:pfv} \textsc{neg} \textsc{3sg.s:opt:ø}-laugh:\textsc{opt:neg} \textsc{purp} \\
\glt `He put his hand over her mouth \textbf{so she wouldn't laugh}.' \\ \hfill [UTOLAF18Jun0805:001]
\end{xlista}
\end{exe}

The undesirability is evaluated from the perspective of the subject of the main clause. The girl who utters (\ref{pova}) wants to escape, but the addressee, her evil ``uncle'', wants her to stay. In (\ref{povb}), people are trying to prevent a man whose only son has just drowned in a boat accident from taking his own life by walking out into the lake.

\begin{exe}
\ex \label{pov}
\begin{xlista}
\ex \label{pova}
\gll ``Sh-taay, tl'uuł  sh-tl'ädn hi-'įh-shüh \textbf{sta-tih-haał} \textbf{ch'a'},'' {\ldots} \textnormal{(S})\\
\textsc{1sg.psr}-paternal.uncle rope \textsc{1sg.psr}-waist \textsc{pp-ø.ipfv:2sg.s:h}-tie:\textsc{ipfv} away-\textsc{inc:aa.ipfv:1sg.s:ø-sg}.go:\textsc{fut:nmz} \textsc{avert} \\
\glt {`}``My uncle, tie a rope around my waist \textbf{lest I run away},'' \ldots'  \\ \hfill \citep[26]{TyoneM1996}
\ex\label{povb}
\gll Ay du' nahugn diniign thüh tl'uul' heh-tsįį ay hii-'i-nih-tl'ǫǫ eł nóó ch'aa-'i-mbeek \textbf{ti-taa-haal} \textbf{ch'a'}. \textnormal{(T)} \\
and \textsc{ct} there moose skin rope:\textsc{poss} \textsc{3pl.s:dh.pfv:h}-make.\textsc{sg.o:pfv} and \textsc{3pl>3sg-pp-n.pfv:h}-tie:\textsc{pfv} and ahead:\textsc{all}  out:\textsc{prmb-3sg.s:ø.ipfv:ø}-swim:\textsc{ipfv:cust} underwater-\textsc{3sg.s:inc:aa.ipfv:ø-sg}.go:\textsc{fut:nmz} \textsc{avert} \\
\glt `And they made a moose skin rope and tied it around him [when] he walked out into the water, \textbf{so he would not go under [and drown himself]}.' \\ \hfill [UTOLAF10Jul0801:023]
\end{xlista}
\end{exe}


\subsubsection{Semantic restrictions on avertive clauses}
\label{subsubsection:avertive-restrictions}

The most important semantic restriction on avertive clauses in \ili{Upper Tanana} is that the event expressed in the avertive clause actually can be prevented by the action in the main clause; i.e.\  the subject of the main clause has control over the event in the avertive clause.\footnote{This contrasts sharply with the situation described by \citet{chapters/wiemer} for several \ili{Slavic} languages, where uncontrollability, i.e.\  the fact that ``neither a discourse participant (in particular the addressee) nor a third party can influence the course of events'', is an integral component of what Wiemer terms ``apprehensional strategies''. This is illustrated abundantly by the examples brought by Wiemer, where possible events to be avoided or prevented, include weather events such as rain or droughts, catching cold or other diseases, or war breaking out.} One can prevent a person from running away or drowning themselves by physically restraining them (\ref{pov}), and screaming or giggling can be prevented, or at least subdued, by placing a hand over a person's mouth (\ref{intrinsic}), but, as the elicited example in (\ref{watch}) shows, one cannot prevent a child from grabbing a knife by merely watching him or her, and thus an avertive linking cannot be used.

\ea
\label{watch}
\gll Ts'iiniin k'à-nįh-tà', \textbf{seeyh} {\textbf{uu-taa-nik}! \textnormal{(N)}} \\
child care-\textsc{qual:ø.pfv:2sg.s:h}-take.care:\textsc{ipfv} knife \textsc{3sg:con-inc:aa.ipfv:ø}-grab:\textsc{fut} \\
\glt `Take care of/watch the baby, he might grab the knife!' \\ \hfill [UTOLAF18Jun0803:003]
\z 

The scenario in the elicited example (\ref{punch}) involves the speaker getting into an argument with an obnoxious stranger at the grocery store. The cue used was `Shut up so I don't beat you up!'. A causal linking is used instead of an avertive one.

\ea \label{punch}
\gll Da-dhį̀į̀-ną́y, \textbf{n-nak-gòt} {\textbf{xah}! \textnormal{(N)}} \\
\textsc{lex-dh.ipfv:2sg.s:ø}-say:\textsc{ipfv:neg} \textsc{2sg.o-inc:qual:aa.ipfv:1sg.s:h}-punch:\textsc{fut} because \\
\glt `Shut up, because I might punch you!' \hfill [UTOLAF18Jun0803:020]
\z

The speakers rejected my suggestion of \textit{nnakgòt ch'à'} `lest I punch you' with the rationale that your shutting up may not be sufficient to keep me from punching you: a rude gesture or facial expression might, after all, provoke me just as much as continued talking would. Another reason for the avoidance of an avertive linking could be that the verb in the main clause is formally negative.\footnote{\textit{Dadhįįnąy} `shut up' is one of a handful of fossilized morphologically formed prohibitives containing the negative prefix \textit{dh}- otherwise lost in \ili{Upper Tanana} \citep{KariJ1993, vanGelderenE2008dene}; \citet[472]{LovickO2020}. The structure of \textit{dadhįįnay} `shut up' is discussed in some more detail in  \citet{LovickO2020ijal}. } I am however uncertain whether speakers recognize this as a negative form due to the degree of fossilization and the lack of negative particles.

In `in-case' scenarios like the elicited example in (\ref{in-case}), avertive clauses can never be used; see \sectref{section:speaking-carefully} for more discussion.

\ea \label{in-case}
\gll  Eh Bill: Tay' t'eey tah-chaan ay xah ch'ale  u-shyii {chąą taathüh} di-dhih-dlay tthiishyoh eh tl'aan dii na-ti-dhag-niign tat'eey u-shyii {dhih-dlah. \textnormal{(T)}} \\
and Bill again \textsc{advzr} \textsc{3sg.s:inc:aa.ipfv:h}-rain:\textsc{fut:nmz} \textsc{3sg} because \textsc{foc} \textsc{3sg.po}-in umbrella \textsc{qual-dh.pfv:1sg.s:ø}-keep.\textsc{pl.o:ipfv:nmz} hat and and what \textsc{prmb-inc-dh.pfv:1sg.s:l}-wear:\textsc{ipfv:cust:nmz} also \textsc{3sg.po}-in \textsc{dh.pfv:1sg.s:ø}-keep.\textsc{pl.o:ipfv} \\
\glt `And Bill: It might rain again, that's why I keep the umbrella in [the backpack] and the hat and my clothes I also keep in it.' \\ \hfill [UTCBAF13Nov1510:035]
\z


\subsection{Alternative constructions}
\label{subsection:alternative-constru}

A number of strategies can be used when an avertive linking is not possible due to semantic or structural restrictions. As shown in the preceding sections, both juxtaposition (\ref{negmaina}) and causal clauses (\ref{negmainb}), (\ref{punch}) and (\ref{in-case}) occur with some frequency, but there are additional options. In `in-case' scenarios, temporal linkings with \textit{ttheh} `before' are a common strategy (\ref{ttheh1}). 

\ea
\label{ttheh1}
\gll ``\textbf{Nah-ógn} \textbf{tah-shyǜǜdn} \textbf{ttheh} na-tsut-dèèł,'' {he-nih. \textnormal{(N)}} \\
\textsc{med}-out:\textsc{ar} \textsc{inc:aa.ipfv:h}-snow:\textsc{fut:nmz} before \textsc{it-1pl.s:inc:opt:d-pl}.go:\textsc{opt} \textsc{3pl.s:ø.ipfv:ø}-say:\textsc{ipfv} \\
\glt {`}``Let's go back before it snows,'' they said.' \hfill [UTCBAF15May0403:039]
\z

\noindent In the speech of at least one speaker, \textit{ttheh} has been extended to have avertive meaning as well (\ref{ttheh2}).

\ea
\label{ttheh2}
\gll \textbf{Ni-h-teedak} \textbf{ttheh} ``Nah-'óg na-h-tet-dag iin,'' nih {hu-'i-chi-jel-nih. \textnormal{(T)}} \\
up-\textsc{3pl.s-inc:dh.pfv:ø-pl}.go:\textsc{pfv} before \textsc{med}-out:\textsc{all} \textsc{prmb-3pl.s-inc:ø.ipfv:d-pl}.go:\textsc{ipfv:cust:nmz} \textsc{pl} say \textsc{3pl.po-pp-lie-indf:3sg.s:qual:ø.ipfv:l}-lie:\textsc{ipfv} \\
\glt `\textbf{To prevent them getting up,} ``There's people walking about outside,'' she said, she lied to them.' \hfill [UTOLVDN13May2909:067]
\z

\noindent (\ref{ttheh2}) is taken from a narrative where a young woman has been abducted by warriors. Her brothers have come to rescue her and are killing her captors outside the tent where she is being held. The men inside the tent hear the noises of the killing, but to prevent them from discovering the slaughter outside, she lies to them, pretending that the noises outside are normal and harmless. Temporal linkings with an avertive interpretation also occur in elicitation (\ref{ttheh3}):

\ea
\label{ttheh3}
\gll U-ch'a' ch'i-nag-'įį \textbf{sh-oo-dah-'aal} \textbf{ttheh}. \textnormal{(T)} \\
\textsc{3sg.po}-from \textsc{indf-n.pfv:1sg.s:l}-hide.self:\textsc{pfv} \textsc{3sg.s>1sg.o-con-inc:qual:aa.ipfv:h}-find:\textsc{co:fut:nmz} before \\
\glt `I'm hiding from him \textbf{to prevent him from finding me.}' \\ \hfill [UTOLAF18Jun0201:016]
\z 

It seems however that the undesirability of the event in the subordinate clause is implicated rather than part of the semantic meaning of the construction. The snow in (\ref{ttheh1}), for example, is not necessarily undesirable, although it is preferable to travel in dry conditions; nor is being found in a game of hide-and-seek in (\ref{ttheh3}) generally problematic. 

In the discussion of the semantics of avertive linkings, the potential existence of a negative purposive linking such as (\ref{intrinsicc}) was noted. Such linkings are attested primarily in elicited examples such as that in (\ref{nolest2}); to facilitate structural comparison, an affirmative purposive clause is provided in (\ref{purpose}).

\ea
\label{nolest2}
\gll ? Olga da-dégn' hǫǫ-heeyh \textbf{k'àt'eey} \textbf{uh-ts'íígn} \textbf{xah}. \textnormal{(N)} \\
{} Olga \textsc{prox}-up:\textsc{all} \textsc{3sg.s:qual:qual:ø.ipfv:ø}-talk:\textsc{ipfv:nmz} \textsc{neg} \textsc{3sg.s:opt:h}-be.sick:\textsc{opt:neg} \textsc{purp} \\
\glt `Olga prays \textbf{so she won't get sick.}' \hfill [UTOLAF14Nov2203:025]
\z 

\ea
\label{purpose}
\gll Olga da-dégn' hǫǫ-heeyh \textbf{na-'ut-shye'} \textbf{xah}. \textnormal{(N)} \\
Olga \textsc{prox}-up:\textsc{all} \textsc{3sg.s:qual:qual:ø.ipfv:ø}-talk:\textsc{ipfv:nmz} \textsc{it-3sg.s:opt:d}-heal:\textsc{pass:opt} \textsc{purp} \\
\glt `Olga prays \textbf{so she'll get healed.}' \hfill [UTOLAF14Nov2203:024]
\z 

\noindent Comparison of (\ref{nolest2}) and (\ref{purpose}) suggests that standard negation of the verb in a purposive clause results transparently in a negated purposive clause. I suspect however that negated purposive clauses such as (\ref{intrinsicc}) or (\ref{nolest2}) are due to influence from \ili{English}. (\ref{nolest2}) was elicited immediately after (\ref{purpose}), using the cue `so she won't get sick'. Both the temporal proximity of the two examples (only one minute elapses between them on the recording) and the \ili{English} cue may have primed the speaker towards producing an \ili{English} calque. Looking through other elicited instances of negated purposive linkings, all of them occur following cues containing `so (that) {\ldots} not'. Even more intriguing is the only clear textual example of this construction, given in (\ref{nolest}).  

\ea
\label{nolest}
\gll \textbf{Ch'its'ą̀ą̀'} \textbf{tsaa-nį́l} \textbf{xah} hǫ̀ǫ̀' {nił-ts'et-niik. \textnormal{(N)}} \\
wrong \textsc{1pl.s:inc:aa.ipfv:ø}-happen:\textsc{fut:neg} \textsc{purp} thus \textsc{recp.po-1pl.s:ø.ipfv:d}-say:\textsc{ipfv:cust} \\
\glt `\textbf{So we won't do it wrong,} we discuss it with each other.' \\ \hfill [UTOLVDN14May0508:114]
\z 

\noindent This utterance occurred about halfway through a 30 minute recording of four elders. For the first ten minutes, one elder talks about how to put on a traditional potlatch, and which pitfalls to avoid. When she is done, a second elder picks up the thread of proper potlatching. (\ref{nolest}) is an utterance produced by the second elder. The subordinate clause is clearly negative despite the absence of the negative particle \textit{k'à(t'eey)}, as evidenced by the tonal pattern. During a transcription and translation session in July 2019 with this same (second) elder, she confirmed that the verb form was formally negative (exaggerating the high tone) and suggested adding the negative particle.\footnote{The negative particle is almost always present, but there are occasional instances without it, as discussed in \citet{LovickOinprep}.} I also asked her to evaluate the avertive clause \textit{ch'its'ą̀ą̀' ts'aanį̀l ch'à'} `lest we do it wrong' (without the tonal pattern associated with standard negation) instead of the negated purposive clause in (\ref{nolest}). The speaker responded that the avertive clause would be vastly preferable to both the original version and the version including \textit{k'àt'eey}, but that it sounded like a rare construction to her.

\largerpage
It is intriguing to compare (\ref{nolest}) with  (\ref{avis-b}), an utterance from earlier in the recording, produced by the first elder.\footnote{Both utterances are uttered in similar contexts: the need to consult with one's close relatives to ensure that gifts at the potlatch are given appropriately. Proper gift-giving at potlatch is a very involved process; see \citet{SimeoneW1995}.}  This structure clearly contains a causal clause with the subordinator \textit{xah} `because', which always triggers nominalization of the subordinate verb. Crucially, the form \textit{hihtaanį̀l} `they might do/happen' is \textit{not} formally negative as evidenced by the low stem tone: \textit{xah} is produced at slightly higher pitch than the stem \textit{nìl}, which bears a lexical low tone.

\ea
\label{avis-b}
\gll Hǫ̀ǫ̀' hudahnih, \textbf{ayaa} \textbf{nts'ą̈̀' t'eey} \textbf{hihtaanį̀l} {\textbf{xah}. \textnormal{(N)}} \\
thus \textsc{3pl.po-qual:2pl.s:ø}-say:\textsc{ipfv} bad \textsc{advzr} \textsc{3pl.s:inc:aa.ipfv:ø}-happen:\textsc{fut:nmz} because \\
\glt `Tell (\textsc{pl}) them, \textbf{because they might do it badly.}' \\ \hfill [UTOLVDN14May0508:053]
\z

After perusing the corpus carefully multiple times, I tentatively conclude that utterances such as (\ref{intrinsicc}), (\ref{nolest2}), or (\ref{nolest}) are \ili{English} calques.


\subsection{Beyond Upper Tanana}
\label{subsection:avert-beyond}

Avertive linkings seem to be attested all across \ili{Northern Dene} (\ref{denaina})--(\ref{slave}),\footnote{Due to my insufficient familiarity with all of these languages, I only provide word-level glosses containing participant and mode/aspect information in this section.} with subordinators that are clearly cognate to \ili{Upper Tanana} \textit{ch'à'}. The constructions in the two Alaskan languages \ili{Dena'ina} (\ref{denaina}) and \ili{Ahtna} (\ref{Ahtna}) are identical to that in Upper Tanana.

\ea \label{denaina}
\gll  Q'uch'a yadi ninya qyeghusłtish dghu qyeł dnish \textbf{tutnitutq'ashi} \textbf{k'u}. \textnormal{(\ili{Dena'ina})} \\
how whatever animal \textsc{3pl>3sg}:skin:\textsc{ipfv:cust} when \textsc{3pl>3sg}:with \textsc{3pl.s}:know:\textsc{ipfv:cust} \textsc{3pl.s}:cut.skin:\textsc{fut:nmz} \textsc{avert} \\
\glt `They knew how to skin any animal \textbf{so as not to damage the skin.}' \\ 
\hfill \citep[260--261]{KalifornskyP1991}
\z

\ea \label{Ahtna}
\gll  Naltsiin 'iin du' xuts'enekaey stanxuhneł'iin \textbf{xuhnał'iił} \textbf{c'a'},  {\ldots} \textnormal{(\ili{Ahtna})} \\
Naltsiin \textsc{pl} \textsc{ct} \textsc{3pl.psr}:children \textsc{3pl>3pl}:hide.away:\textsc{pfv} \textsc{3pl>3pl}:see:\textsc{fut} \textsc{avert} \\
\glt `The Naltsiin had hidden their children \textbf{so [the Tailed People] wouldn't see them,}  \ldots' \hfill \citep[55]{KariJ1986}
\z


 In the two Canadian languages, the construction looks a little different. This is partially due to differences between the language groups; while future and negative polarity are inflectional categories in the Alaskan languages, they are expressed by postverbal particles or clitics in the Canadian ones. This likely explains the imperfective form in the avertive clause in the \ili{Dëne Sųłiné} example (\ref{denesuline}). The \ili{Slave} example in (\ref{slave}) uses an optative form in the avertive clause, but most surprising is the fact that this form contains the negative enclitic =\textit{híle} (this enclitic is partly absorbed into the stem). The same is true for all other avertive linkings cited by \citet{RiceK1989}, suggesting that the construction in \ili{Slave} has undergone some changes, possibly under the influence of \ili{English} contact.

\ea \label{denesuline}
\gll Hustón=ë \textbf{yëkéth} \textbf{ch'á.}  \textnormal{(\ili{Dëne Sųłiné})}\\
\textsc{1sg>3sg}:hold:\textsc{ipfv=fut} \textsc{3sg}:fall.over:\textsc{ipfv} \textsc{avert} \\
\glt `I'll hold [the swingset] \textbf{so it doesn't fall over.}' \\ \hfill [deslas-ANR-2016-07-10-C] 
\z

\ea \label{slave}
\gll   \textbf{Daniel} \textbf{yegúhʔále} \textbf{ch'á} goghádehk'a. \textnormal{(\ili{Slave})} \\
Daniel \textsc{3sg>3sg}:find.\textsc{co:opt:neg} \textsc{avert} \textsc{1sg>3sg}:throw.\textsc{co:pfv} \\
\glt `I threw it \textbf{so Daniel wouldn't find it.}' \hfill \citep[1262]{RiceK1989} 
\z


\subsection{Origin and other uses of \textit{ch'à'}}
\label{subsection:avertive-origin}
\largerpage


The origin of the avertive subordinator is likely the Proto-\ili{Dene} postposition \textsc{po}+\textit{*k'ʸaˑn̠}' `away, apart from \textsc{po}' \citep{LeerJ1996CA}.\footnote{A similar origin is observed by \citet{chapters/francois} for a comparable subordinator in several Torres-Banks languages.} In many daughter languages, cognates occur as a free or bound postposition with meanings such as `separating \textsc{po}, emanating from \textsc{po}', or adverbial-derivational verb prefixes with the meaning `out and away'. In the Alaskan languages (apparently not beyond), the meanings `warding off \textsc{po}, protecting against/from \textsc{po}, keeping \textsc{po} away' are also common.

In Upper Tanana, cognates of the Proto-\ili{Dene} morpheme include the bound adverbial-derivational morpheme \textit{ch'a}- `out and away' which is restricted to motion verbs (\ref{adv-der}). This morpheme has purely spatial meaning without connotations of avoidance or fear.

\begin{exe}
\ex\label{adv-der}
\begin{xlista}
\ex 
\gll {Ch'a-y-dih-tįį {\ldots} \textnormal{(T)}} \\
out-\textsc{3sg>3sg-qual-aa.pfv:h}-handle.\textsc{ao:pfv} \\
\glt `She took him out [from between two spruce trees] \ldots' \\ \hfill [UTOLAF09Jun2202:015]
\ex
\gll Nóó ch'a-ts'i-nįį-kįį eh {\ldots} \textnormal{(T)}\\
ahead:\textsc{all} out-\textsc{1pl.s-n.pfv:ø}-go.by.boat:\textsc{pfv} and\\
\glt `We start out by boat and \ldots' \hfill [UTOLAF09Jun2902:011]
\end{xlista}
\end{exe}

The free postposition \textsc{po}+\textit{ch'à'} `away from \textsc{po}', by contrast, does have strong connotations of avoidance (\ref{avoida}, \ref{avoidb}), hiding (\ref{avoidc}), or escape (\ref{avoidd}). This last example was a prompt by me that was flagged as problematic by the speakers, who suggested (\ref{avoide}) instead. A sentence like (\ref{avoidd}) would be appropriate if the postpositional object were some sort of evil creature, but not if it is a respected elder like Sherry—in that situation, the (purely spatial) ablative postposition like \textsc{po}+\textit{ts'änh} `from \textsc{po}' in (\ref{avoide}) has to be used. 

\begin{exe}
\ex 
\label{avoid}
\begin{xlista}
\ex\label{avoida}
\gll Nts'ą' na-tthi-ts'el-dak shyįį' t'eey, nih, \textbf{huk'an} \textbf{ch'a'}. \textnormal{(T)} \\
and \textsc{cont}-out-\textsc{1pl.s:ø.ipfv:l-pl}.go:\textsc{ipfv:cust} only \textsc{advzr} say:\textsc{ipfv} burn-out from \\
\glt `And we were just running, she said, \textbf{from that burn-out}.' \\ \hfill [UTOLAF09Jun2402:028] 
\ex \label{avoidb}
\gll \textbf{Hu-ch'à'} {i-tal-ch'ee {\ldots} \textnormal{(N)}} \\
\textsc{3pl.po}-from \textsc{?-3sg.s:inc:aa.ipfv:l}-be.aloof:\textsc{fut} \\
\glt [Money, success] will stay \textbf{away from them} \ldots' \\ \hfill [UTOLAF12Jul1203:067]
\ex \label{avoidc}
\gll S-seeyy' sta-tnįh-'įįl {\textbf{sh-ch'a'} {\ldots} \textnormal{(T)}} \\
\textsc{1sg.psr}-knife:\textsc{poss} away-\textsc{fut:qual:aa.ipfv:2sg.s:h}-hide:\textsc{fut:nmz} \textsc{1sg.po}-from \\
\glt `You (\textsc{sg}) must have hidden my knife \textbf{away from me} \ldots' \\ \hfill [UTOLVDN13May2804:077]
\ex \label{avoidd}
\gll ? \textbf{Sherry} \textbf{shyah} \textbf{ch'à'} {ih-haał. \textnormal{(N)}} \\
{} Sherry house from \textsc{aa.ipfv:1sg.s:ø-sg}.go:\textsc{ipfv:prog} \\
\glt (Intended: `I'm coming \textbf{from Sherry's house}') \\ Comment: `Sounds like you're getting away from her.' \\ \hfill [Prompt by OL for UTOLAF18Jun0101:087]\\
\ex \label{avoide}
\gll \textbf{Sherry} \textbf{shyah} \textbf{ts'änh} {ih-haał. \textnormal{(N)}} \\
Sherry house from \textsc{aa.ipfv:1sg.s:ø-sg}.go:\textsc{ipfv:prog} \\
\glt `I'm coming \textbf{from Sherry's house}.' \hfill [UTOLAF18Jun0101:087]
\end{xlista}
\end{exe}

\textsc{po}+\textit{ch'à'} also has non-spatial meanings as demonstrated in (\ref{cha}), but it retains its negative connotations. 

\begin{exe}
\ex 
\label{cha}
\begin{xlista}
\ex \label{chaa}
\gll Chih nah-'ógn ni-'eeł-ts'e-łeek {nts'ą' t'eey} t'axoh ti-ts'at-ndak {\textbf{u-ch'a'}. \textnormal{(T)}} \\
also \textsc{med}-out:\textsc{ar} \textsc{term}-trap-\textsc{1pl.s:ø.ipfv:l}-handle.\textsc{pl.o:ipfv:cust} \textsc{advzr} finally \textsc{lex-1pl.s:aa.pfv:d}-be.tired:\textsc{ipfv} \textsc{3sg.po}-from \\
\glt `We would be setting traps too out there; finally we'd get tired \textbf{of it}.' \\ \hfill [UTCBVDN13Nov1501:096]
\ex \label{chab}
\gll \textbf{Ch'indziidn} \textbf{ch'à'} {dhih-säł. \textnormal{(N)}} \\
bee from \textsc{dh.pfv:1sg.s:ø}-scream:\textsc{pfv} \\
\glt `I screamed \textbf{for fear of the bees}.' \hfill [UTOLAF18Jun0805:037]
\end{xlista}
\end{exe}

\noindent (\ref{chab}) comes from elicitation; the scenario given to speakers was that I was stung twelve times by bees the last time I stole honey: now the bees are after me again and I scream with fear.

The textual example (\ref{avert-pp}) shows that the postposition \textsc{po}+\textit{ch'à'} can also be used with clearly avertive function. The speaker explains why silence is expected during condolence visits. (\ref{avert-ppa}) paints a scenario where the visitor is not silent; the free pronoun \textit{ay} in (\ref{avert-ppb}) refers to that scenario. (\citet[90–92]{BassoK1990} reports a similar injunction against unnecessary conversation in encounters with the recently bereaved in the \ili{Western Apache} context.)

\begin{exe}
\ex 
\label{avert-pp}
\begin{xlista}
\ex \label{avert-ppa}
\gll N-èh stoo-tnèl-nay èh dįį-nįįł nts'ą̈̀', shyą̀ą̀' t'eey ha-tha-dįh-nay dè' ch'its'ą̀ą̀' tah {dįį-nįįł. \textnormal{(N)}} \\
\textsc{2sg.po}-with away:\textsc{ar.s-inc:qual:dh.pfv:l-co}.moves:\textsc{pfv:nmz} with \textsc{inc:qual:aa.ipfv:2sg.s:ø}-say:\textsc{fut} and for.nothing \textsc{advzr} out-mouth-\textsc{qual:aa.pfv:2sg.s:h}-move.\textsc{co}.carelessly\textsc{:pfv:nmz} if wrong among \textsc{inc:qual:ø.ipfv:2sg.s:ø}-say:\textsc{fut} \\
\glt `You (\textsc{sg}) will talk when [your mind] is not all the way there, and if you (\textsc{sg}) just blab for no reason, you (\textsc{sg}) will say something wrong.'\footnote{This sentence contains two fascinating morphological idioms. \textit{Neh stootnelnay eh} `when your mind is not all the way there' can be glossed more literally as `when a compact object has moved independently away, affecting you', while \textit{hathadįhnay} `you just blab' could be glossed as `you move a compact object carelessly out using your mouth'.} \\ \hfill [UTOLVDN14May0508:013]
\ex\label{avert-ppb}
\gll \textbf{Ay} \textbf{ch'à'} shyą̀ą̀' {dhįį-dah. \textnormal{(N)}} \\
\textsc{3sg} \textsc{avert} for.nothing \textsc{dh.pfv:2sg.s:ø-sg}.sit:\textsc{ipfv} \\
\glt `\textbf{To avoid that [i.e., saying the wrong thing],} sit (\textsc{sg}) for nothing [without talking]. \hfill [UTOLVDN14May0508:014]
\end{xlista}
\end{exe}


Although the postposition \textsc{po}+\textit{ch'à'} may have connotations of fear, the subordinator does not. Unlike what \citet[296–297]{LichtenberkF1995} describes for the Austronesian language \ili{To'aba'ita}, \textit{ch'à'} is not attested in the corpus as a subordinator when the verb of the main clause is a predicate of fearing; instead, the complementizer \textit{t'eey} is used (\ref{sleep}):

\ea
\label{sleep}
\gll \textbf{Tee-te'} \textbf{t'eey} {neljit! \textnormal{(T)}} \\
\textsc{3sg.s:inc:dh.pfv:ø-sg}.sleep:\textsc{pfv} \textsc{comp} \textsc{3sg.s:qual:ø.pfv:l}-be.scared:\textsc{ipfv} \\
\glt `She was scared \textbf{to fall asleep}!' \hfill [UTCBAF13Nov0501:086]
\z

\textit{Ch'à'} could also not be elicited with this function as shown in (\ref{afraid}); the cues were `I'm afraid the dog might bite me', `I'm afraid they might kill us', and `I'm afraid he will find me'. 

\begin{exe}
\ex\label{afraid}
\begin{xlista}
\ex
\gll Łįį sh-taa-'aał {m-i-nag-jit. \textnormal{(T)}} \\
dog \textsc{1sg.o-3sg.s:inc:aa.ipfv:ø}-bite:\textsc{fut} \textsc{3sg.po-pp-qual:ø.pfv:1sg.s:l}-be.afraid:\textsc{ipfv} \\
\glt `The dog might bite me, I'm afraid of it.' \hfill [UTOLAF18Jun0201:001]\\
\ex 
\gll Han iin hu-'i-nagn-jit, nee-h-taa-xąą k'eh {hel-tsįį. \textnormal{(N)}}\\
\textsc{dem} \textsc{pl} \textsc{3pl.po-pp-qual:ø.pfv:1sg.s:l}-be.scared:\textsc{ipfv} \textsc{1pl.o-3pl.s-inc:aa.ipfv:ø}-kill.\textsc{pl.o:fut} like \textsc{3pl.s:dh.pfv:l}-seem:\textsc{ipfv}\\
\glt `I'm afraid of them, they look like they might kill us.' \\ \hfill [UTOLAF18Jun0101:008]\\
\ex
\gll Idzii {sh-oo-dah-'aał. \textnormal{(T)}} \\
scary \textsc{1sg.o-3sg.s:qual-inc:qual:aa.ipfv:h}-find.\textsc{co:fut} \\
\glt `Scary, he will find me.' \hfill [UTOLAF18Jun0201:015]\\
\end{xlista}
\end{exe}


\subsection{Avertive clauses and speech acts}
\label{subsubsection:avertive-speech}

\citet[24]{DixonRMW2009} notes that typologically, the main clause in an avertive construction (``possible consequence linking'' in his terminology) is ``most often a positive imperative  {\ldots} or a negative one''.\footnote{The basis for his claim is unknown to me.} In Upper Tanana, this is not the case. As mentioned above, avertive clauses cannot be used at all with negative requests. With affirmative requests, they are rare. \tabref{tab:avertive-freq} shows that almost all naturalistic instances of avertive linking occur in declarative speech acts.

\begin{table}
\caption{Speech act types in avertive linkings}
\label{tab:avertive-freq}
 \begin{tabular}{rrrrr}
  \lsptoprule
  declarative & question  & request & hortative & total \\   \midrule
  14 & 2 & 1 & 1 & 18\\
  \lspbottomrule
 \end{tabular}
\end{table}

The total number of naturalistic examples of avertive clauses is so small (18 instances in about 10,000 utterances) that the relative frequencies have to be taken with a grain of salt. Yet it does appear as though Dixon's claim about the predominance of avertive linkings in requests cannot be maintained for Upper Tanana. Since avertive linkings in requests can however be elicited without difficulty, the rarity of this construction is likely to have pragmatic rather than grammatical reasons. I suggest here that this is due to the availability of an urgent request construction that points to, but does not spell out, potential undesirable consequences of non-compliance. These urgent requests are described in \sectref{section:apprehensional}.

\section{Urgent requests}
\label{section:apprehensional}

The structure of urgent requests is described in \sectref{subsection:urgent-structure}. This is followed by a discussion of the origin of this construction in \sectref{subsection:urgent-origin}, and of its semantics in \sectref{subsection:urgent-semantics}.

\subsection{Structure}
\label{subsection:urgent-structure}

Urgent requests are in the optative and always contain the particle  \textit{dè'} (\ref{opt-de}). 

\begin{exe}
\ex \label{opt-de}
\gll Jin dǫh-'à' dè'. \textnormal{(B)}\\
this \textsc{qual:opt:2sg.s:vv}-sing:\textsc{opt} please? \\
\glt `Sing (\textsc{sg}) this [song], you (\textsc{sg}) should sing this [song], please sing (\textsc{sg}) this [song].' \hfill [UTOLAF10Jul0202:196]
\end{exe}

\noindent In discussing (\ref{opt-de}) and others like it, \citet[269]{LovickO2016} suggests that the translation `please' might be appropriate for this particle and that it serves to mitigate the request. More recent discussions with speakers, however, suggest a different function for this particle: that of adding urgency to a request, resulting in the translation `make sure you sing this [song]' for (\ref{opt-de}), as discussed in \sectref{subsection:urgent-semantics}.

\subsection{Origin}
\label{subsection:urgent-origin}

The origin of urgent requests with \textit{dè'} is likely a conditional request, which are a special type of conditional linking. In declarative conditional linkings, the conditional clause is marked by \textit{dè'} `if, when.\textsc{fut}'. (This polysemy is common in \ili{Dene} languages, see also \citealt{AnismanA2019}.) The consequent clause is usually in the future.

\ea \label{conditional}
\gll Ay \textit{and} ay tl'ààn udzih dhèh-xįį dè' nah-'ǫ́ǫ́ nìì-taa-łeeł. \textnormal{(N)} \\
and and and then caribou \textsc{3sg.s:qual:dh.pfv:h}-kill.\textsc{sg.o:pfv} if \textsc{med}-out:\textsc{abl} \textsc{term:it-3sg.s:inc:aa.ipfv:ø}-handle.\textsc{pl.o:pfv} \\
\glt `And if he'd kill a caribou, he'd bring it from out there.' \\ \hfill [UTOLVDN07Nov2901:020]
\z

Conditional requests differ from conditional declaratives in two respects: the consequent is usually in the optative, and both clauses contain the particle \textit{dè'}~(\ref{cond-de}).

\begin{exe}
\ex \label{cond-de}
\begin{xlista}
\ex \label{cond-dea}
\gll ``K'àt'eey k'àt'eey hǫǫsú' m-i-hǫl-náy \textbf{dè'} t'eey, \textbf{at-naah} \textbf{dè'},'' nih. \textnormal{(N)} \\
\textsc{neg} \textsc{neg} good \textsc{3sg.p-pp-ar:ø.pfv:l}-taste:\textsc{ipfv:nmz} if even \textsc{opt:2pl.s:d}-drink:\textsc{opt} \textsc{ur} say:\textsc{ipfv} \\
\glt {`}``Even if it doesn't taste good, make sure you (\textsc{pl}) drink it,'' he said.' \\ \hfill [UTOLAF18Jun0401:013]
\ex\label{cond-deb}
\gll Ch'i-dhih-'eeł \textbf{de'}, \textbf{k'a'} \textbf{ǫǫ-keet} \textbf{de'}, \textbf{huugn} \textbf{ta-'ǫl-shyeek} \textbf{de'},  {\ldots} \textnormal{(T)} \\
\textsc{indf.o-dh.pfv:2sg.s:h}-trap:\textsc{pfv} if gun \textsc{opt:2sg.s:ø}-buy:\textsc{opt} \textsc{ur} etc. distribute-\textsc{opt:2sg.s:l}-handle.\textsc{pl.o:opt:dist} \textsc{ur} \\
\glt `And if you (\textsc{sg}) trap something, make sure you (\textsc{sg}) buy guns with it, make sure you (\textsc{sg}) distribute it [at potlatch] \ldots' \\ \hfill [UTCBAF13Nov0401:070]
\end{xlista}
\end{exe}

\noindent The presence of \textit{dè'} `if, when.\textsc{fut}' in both clauses in conditional requests suggests that the consequent clause sets up a further condition with an unspoken consequence: even if the broth doesn't taste good, make sure you drink it; if you do not drink it, you will  die of starvation.\footnote{The broth in question is prepared from moose droppings and consumed only under the most dire circumstances.} If you have good luck trapping, make sure to distribute your wealth through a potlatch; if you fail to do so, you will no longer be successful when trapping. Full conditional requests like those in (\ref{cond-de}) thus are instances of insubordination \citep{evans2007}, where the main clause is unspoken, and where the (second) conditional clause is modified by the first conditional clause. Urgent requests such as (\ref{opt-de}) lack the first conditional clause. 

\subsection{Semantics}
\largerpage
\label{subsection:urgent-semantics}

Urgent requests with \textit{dè'} imply the presence of undesirable consequences of non-compliance, such as snow out of season in the case of (\ref{opt-de}), starvation in (\ref{cond-dea}), or poor luck in trapping in (\ref{cond-deb}). More examples are given in (\ref{more-de}). 

\begin{exe}
\ex \label{more-de}
\begin{xlista}
\ex\label{more-dea}
\gll Tl'aan hu-na-chin-'ah-deeł tl'aan hu-b-e-kon'-nah-shyiił de' hu-'ah-xąą de', {nee-he-nih. \textnormal{(T)}} \\
and \textsc{3pl.p}-surround-group-\textsc{opt:2pl.s:ø-pl}.move:\textsc{opt} and \textsc{3pl.p-?-pp}-fire-\textsc{qual:opt:2pl.s:h}-strike.\textsc{lro:opt} \textsc{ur} \textsc{3pl.o-opt:2pl.s:h}-kill.\textsc{pl.o:opt} \textsc{ur} \textsc{1pl.p-3pl.s:ø.ipfv:ø}-say:\textsc{ipfv}\\
\glt `Then gang (\textsc{pl}) up on them and make sure you (\textsc{pl}) set fire to them and kill them, they said to us.' \hfill [UTCBVDN14Jul1801:127] \\
\ex\label{more-deb}
\gll Nts'ą̀ą̀' n-dih-nih niign t'eey dǫǫ-dį̀' dè', {u-dih-nih. \textnormal{(N)}}\\
whatever \textsc{2sg.p-qual:1sg.s:ø}-say:\textsc{ipfv} \textsc{comp} \textsc{advzr} \textsc{qual:opt:2sg.s:ø}-do:\textsc{ipfv} \textsc{appr} \textsc{3sg.p-qual:1sg.s:ø}-say:\textsc{ipfv} \\
\glt `Make sure you (\textsc{sg}) do [exactly] whatever I tell you (\textsc{sg}), I said to him.' \hfill [UTOLVDN13May2004:067]
\end{xlista}
\end{exe}

\noindent (\ref{more-dea}) is taken from a text about World War II, where US army representatives instructed the people of Tetlin village on what to do if they were attacked by the Japanese. Non-compliance here would have resulted not only in the death of many villagers, but also in the Japanese establishing a foothold on the Alaskan mainland, something the US army desperately wanted to avoid. (\ref{more-deb}) (discussed at length in \citealt{LovickO2016}) is from a text where the speaker and her brothers are attacked by a bear, but only have a single shot between them. The speaker is trying to protect her brothers from the bear; when her older brother starts arguing, she shuts him up with (\ref{more-deb}). In this example, the consequence of non-compliance would be being killed by the bear. Generally, \textit{dè'} is used in urgent requests where non-compliance would result in very undesirable consequences.

It should be noted that this modified analysis of \textit{dè'} as an urgency marker does not necessarily invalidate the claim that it mitigates a request, as put forward in \citet[269]{LovickO2016}. As \citet{LovickO2016} notes, reiterating arguments also made by \citet{ScollonRScollonS1981} or \citet{RushforthSChisholmJ1991}, uttering requests is considered to be interfering with another's freedom of action and thus undesirable in and of itself. Using \textit{dè'} in a request draws attention to the undesirable consequences implicit in non-compliance, which in turn provides motivation for uttering the request in the first place -- a form of mitigation. 

Thus requests containing \textit{dè'} serve the speaker's need to point to the necessity of compliance but crucially never spell out the possible consequences of non-compliance, trusting the addressee to be able to work this out for themselves.\footnote{This looks like the reversal of the indirect requests described by \citet{chapters/francois} for \ili{Mwotlap}, where the undesirable consequence is spelled out, but no request is uttered.} Spelling out the desirable consequences of compliance is not compatible with the \ili{Upper Tanana} practice of speaking carefully about the future described in \sectref{subsection:careful-action}.

\subsection{Beyond Upper Tanana}
\label{subsubsection:beyond-urgent}

Similar constructions do not seem to be reported for related languages.


\section{Speaking carefully}
\label{section:speaking-carefully}

In this section, I discuss several aspects of avertive linkings and urgent requests in greater detail, showing how they reflect general cultural patterns in \ili{Alaskan Dene}: the keen awareness of which things are beyond an individual's control, and the need to be prepared for those things that are not (\sectref{subsection:not-afraid}) as well as the principle of non-interference and respecting the competence of others (\sectref{subsection:other-avertive}).

\subsection{Control and preparedness}
\label{subsection:not-afraid}

In \sectref{subsection:avertive-semantics}, I showed that avertive linkings cannot be used in `in-case' scenarios but only when the event in the avertive clause is actually under the control of the subject of the main clause. Speakers proved to be quite sensitive about this, consistently rejecting examples where control was absent or inadequate. As discussed in \sectref{subsection:alternative-constru}, there are several other constructions that can be used under these circumstances.

Yet, precisely because the Upper Tanana are aware of their limited control over their environment, they employ a circumspect mindset requiring them to be prepared for any eventuality. This had some surprising consequences for fieldwork. One of the examples of `in-case' scenarios  comes from a storyboard\footnote{Storyboards are sequences of images serving as stimulus for a speaker to tell the story in their own words. They allow the researcher to control the context of an utterance without having to use the contact language. For discussion of this methodology, see \citet{BurtonSMatthewsonL2015}.} \citep{vanderKlokJ2013}; it is given in (\ref{in-case2}), repeated from (\ref{in-case}) above.

\ea \label{in-case2}
\gll  Eh Bill: Tay' t'eey tah-chaan ay xah ch'ale  u-shyii {chąą taathüh} di-dhih-dlay tthiishyoh eh tl'aan dii na-ti-dhag-niign tat'eey u-shyii {dhih-dlah. \textnormal{(T)}} \\
and Bill again \textsc{advzr} \textsc{3sg.s:inc:aa.ipfv:h}-rain:\textsc{fut:nmz} \textsc{3sg} because \textsc{foc} \textsc{3sg.po}-in umbrella \textsc{qual-dh.pfv:1sg.s:ø}-keep.\textsc{pl.o:ipfv:nmz} hat and and what \textsc{prmb-inc-dh.pfv:1sg.s:l}-wear:\textsc{ipfv:cust:nmz} also \textsc{3sg.po}-in \textsc{dh.pfv:1sg.s:ø}-keep.\textsc{pl.o:ipfv} \\
\glt `And Bill: It might rain again, that's why I keep the umbrella in [the backpack] and the hat and my clothes I also keep in it.' \\ \hfill [UTCBAF13Nov1510:035]
\z

\noindent Bill's behavior -- he arrives at work wearing a coat, a hat, and carrying an umbrella on a fine day -- may strike someone socialized in the mainstream American or Canadian worldview as ridiculous.  \ili{Upper Tanana Dene}, however, consider this perfectly sensible behavior due to the emphasis on \textit{ts'edòònih} discussed in \sectref{subsection:careful-action}. Preparedness as that shown by Bill is valued and in fact encouraged as shown in (\ref{in-case3}), which is taken from a text about \textit{ts'edòònih}:
 
\begin{exe}
\ex \label{in-case3}
\begin{xlista}
\ex 
\gll Nithaad dą̀' sh-nąą iin èh sh-tà' eh ``Hah-'ógn elih tah hǫǫsú', hǫǫsú' din-tl'ǫǫk hah'ógn ts'edòònih xah,'' {nee-he-niik. \textnormal{(N)}}\\
long.time when:\textsc{pst} \textsc{1sg.psr}-mother \textsc{pl} and \textsc{1sg.psr}-father and \textsc{ntrl}-outside:\textsc{ar} cold when well well \textsc{qual:ø.ipfv:2sg.s:d}-dress:\textsc{ipfv:cust} \textsc{ntrl}-outside:\textsc{ar} never.know because \textsc{1pl.po-3pl.s:ø.ipfv:ø}-say:\textsc{ipfv:cust} \\
\glt `A long time ago, our parents used to say to us, ``When it is cold outside, dress (\textsc{sg}) well, because you never know.''' \\ \hfill [UTOLAF19Jan0501:010]
\ex 
\gll ``Elók hǫǫ-łįį nįį-thänh dè' n-ch'ìl' {tįį-łeek.'' \textnormal{(N)}}\\
hot \textsc{ar.s:ø.pfv:ø}-be:\textsc{ipfv} \textsc{qual:ø.pfv:2sg.s:ø}-think:\textsc{ipfv} if \textsc{2sg.psr}-clothes:\textsc{poss} \textsc{inc:ø.ipfv:2sg.s:ø}-handle.\textsc{pl.o:ipfv:cust} \\
\glt {`}``Even if you (\textsc{sg}) think it is hot, take your (\textsc{sg}) clothes.''' \\ \hfill [UTOLAF19Jan0501:011]
\ex 
\gll N-èh n-èh tįį-łeek nah-'ógn nts'ą̀ą̀' u-taa-nį̀ł nts'ą̈̀' hi-ts'-in-níígn {xah. \textnormal{(N)}} \\
\textsc{2sg.po}-with \textsc{2sg.po}-with \textsc{inc:ø.ipfv:2sg.s:ø}-handle.\textsc{pl.o:ipfv:cust} \textsc{med}-outside:\textsc{ar} \textsc{advzr} \textsc{ar.s-inc:aa.ipfv:ø}-happen:\textsc{fut} \textsc{advzr} \textsc{pp-pej-ø.pfv:2sg.s:d}-not.know:\textsc{ipfv} because \\
\glt {`}``Take them with you because you don't know what will happen out there.''' \hfill [UTOLAF19Jan0501:012]
\end{xlista}
\end{exe}

In a way, the cultural emphasis on preparedness could be considered one reason for the rarity of avertive linkings in the narrative corpus. 

\subsection{Respecting the competence of others}
\label{subsection:other-avertive}

Looking at the distribution of avertive linkings and urgent requests by genre (\tabref{tab:avertive-freq2}), it becomes apparent that they are used with very different audiences.

\begin{table}
\caption{Avertive linkings and urgent requests by genre}
\label{tab:avertive-freq2}
 \begin{tabular}{l rr } 
  \lsptoprule
      & avertive & urgent requests  \\   \midrule
past events & 12 & 44 \\
teachings & 6 & 3 \\\midrule
total & 18 & 47\\ 
  \lspbottomrule
 \end{tabular}
\end{table} 

``Past events'' includes the three narrative genres of old-time stories, histories, and personal narratives identified in \citet{LovickO2012b-genre}.  As pointed out by Jones \& Thompson (\citeyear[ix]{JonesEThompsonC1990}) in the context of old-time stories, these stories are told ``to entertain, educate and inspire -- to cause the reader or listener to think'' (see also Bergelson (\citeyear{BergelsonM2019}) on lessons implicit in personal narratives). The typical audience for this type of narrative is composed of community members. Crucially, the lessons contained in these stories are expressed implicitly and require the audience to unpack them, using the cultural knowledge at their disposal. ``Teachings'' refers to a non-traditional genre where speakers explain cultural knowledge to an audience consisting of community outsiders such as linguists and anthropologists, or to young people who did not grow up in the traditional way. The explanations in such texts are very explicit and often contain rationales for the ritualistic acts discussed in \sectref{subsection:careful-action}. It is thus not surprising that a third of all examples of avertive linkings come from teachings: these are explanations of cultural practices aimed at those who are unfamiliar with them ((\ref{undesirablea})--(\ref{undesirableb}) fall into that category). It is also not surprising that there are so few urgent requests in these texts: leaving aside the relative rarity of non-declaratives in this text type, speakers are also cognizant of the fact that the audience of a ``teaching'' text may not be able to interpret the request correctly. Instead, urgent requests occur predominantly in stories where characters well-versed in Upper Tanana culture interact with one another (see \citet{LovickO2016} for a discussion of using stories for the investigation of requests). The traditional audience for these texts does need to be told the consequences of (non-)compliance. Explicitly stating the (desirable) outcome of compliance with the urgent request would be a violation of the practice of careful speaking. Stating the (undesirable) consequence of non-compliance, on the other hand, could suggest that the addressee is not able to think for themselves or ``does not know their \textit{įįjih}'', which is not compatible with the principle of non-interference and respect. In either case, verbalizing the potential consequences in an urgent request would collide with fundamental principles of Upper Tanana culture.

Together, avertive linkings and requests with \textit{dè'} fulfill an \ili{Upper Tanana} speaker's need to remind the addressee of circumspect action, while at the same time treating them as a competent person who ``knows their \textit{įįjih}''. 



\section{Conclusion}
\label{section:conclusion}

In this paper, I described avertive linkings and urgent requests in \ili{Upper Tanana Dene}, and discussed how they reflect the circumspect mindset of Upper Tanana people. 

Avertive linkings are a type of adverbial linking involving the subordinator \textit{ch'à'} `\textsc{avert}, lest'. The avertive clause is usually in the future, more rarely in the inceptive perfective mode. It is always formally affirmative. The main clause has no mode restrictions and also has to be formally affirmative. There are no other formal restrictions on avertive linkings. Semantically, avertive linkings are limited to situations where an undesirable consequence (expressed in the avertive clause) can be avoided by taking the preemptive action expressed in the main clause. This linking is very sensitive to the notion of what can and cannot be controlled by human action and thus cannot be used in `in-case' scenarios. It also cannot be used in linkings where the main predicate is one of fear. 

Somewhat unusually typologically, avertive linkings with a request as main clause are rare in Upper Tanana. I argue in this paper that this is likely due to the fact that requests are problematic in \ili{Northern Dene} in general, and that spelling out the undesirable consequences to be avoided could suggest to the hearer that they are thought to be lacking in cultural knowledge. Instead, there is a dedicated strategy for urgent requests formed using the particle \textit{dè'}. This type of request likely has arisen as a insubordinate conditional structure. The undesirable consequence of non-compliance is implicated, but never spelled out.

The existence and usage patterns of these constructions reflect the circumspect mindset of the Upper Tanana people, i.e., the need to act and speak mindfully. Avertive linkings spell out the undesirable consequences that can be avoided by taking the action in the main clause. They occur primarily in texts directed at community outsiders who often are unaware of the many, often ritualistic, actions performed to ensure a favorable outcome. Requests formed with \textit{dè'} point to the existence of undesirable consequences, but crucially never spell these out: the speaker assumes that the hearer has the required cultural knowledge to do so. Thus, such requests usually are directed towards a member of the community, who can be relied on to have this knowledge. In this fashion, avertive linkings and requests with \textit{dè'} complement each other in the domain of circumspection, while respecting the addressee's cultural competence.


\section*{Abbreviations}

\begin{tabularx}{.5\textwidth}{@{}lQ@{}}
\textsc{aa} & \textit{aa}-conjugation \\
\textsc{abl} & ablative \\
\textsc{advzr} & adverbializer \\
\textsc{all} & allative \\
\textsc{ar} & areal \\
\textsc{avert} & avertive \\
B & Beaver Creek dialect \\
\textsc{comp} & complementizer \\
\textsc{co} & compact object \\
\textsc{cont} & continuative \\
\textsc{ct} & contrastive topic \\
\textsc{cust} & customary \\
\textsc{d} & D-voice/valence maker \\
\textsc{dh} & \textit{dh}-conjugation \\
\textsc{foc} & focus \\
\textsc{fut} & future \\
\textsc{h} & H-voice/valence marker \\
\textsc{imp} & imperative \\
\textsc{ipfv} & imperfective \\
\textsc{inc} & inceptive \\
\textsc{indf} & indefinite \\
\textsc{it} & iterative \\
\textsc{l} & L-voice/valence marker \\
\textsc{lex} & lexical morpheme \\
\textsc{lro} & long rigid object \\
\textsc{med} & medial distance \\
N & Northway dialect \end{tabularx}%
\begin{tabularx}{.5\textwidth}{@{}lQ@{}}
\textsc{n} & \textit{n}-conjugation \\
\textsc{nmz} & nominalizer \\
\textsc{neg} & negative \\
\textsc{ntrl} & neutral distance \\
\textsc{o} & direct object \\
\textsc{opt} & optative \\
\textsc{po} & postpositional object \\
\textsc{pej} & pejorative \\
\textsc{pfv} & perfective \\
\textsc{pl} & plural \\
\textsc{poss} & possessed \\
\textsc{pp} & postposition \\
\textsc{prmb} & perambulative \\
\textsc{prog} & progressive \\
\textsc{proh} & prohibitive \\
\textsc{psr} & possessor \\
\textsc{purp} & purposive \\
\textsc{qual} & qualifier \\
\textsc{recp} & reciprocal \\
\textsc{refl} & reflexive \\
S & Scottie Creek dialect \\
\textsc{s} & subject \\
\textsc{sg} & singular \\
T & Tetlin dialect \\
\textsc{term} & terminative \\
\textsc{ur} & urgent request \\\\
\end{tabularx}


\section*{Acknowledgements}

I want to express my deep gratitude to all of the Upper Tanana speakers who so patiently teach me about their language and their culture. From Northway: Mr. Roy Sam and Mrs. Avis Sam as well as Mrs. Sherry Demit Barnes. From Tetlin: Mr. Roy H. David, Sr., and Mrs. Ida Joe. \textit{Tsinah'įį} for your hard work, and for all the laughter. 

\textit{Marsi choh!} also goes to the Doyon Languages Online project led by Allan Hayton and to Jim Kari and the College of Rural and Cross-Cultural Studies at the University of Alaska Fairbanks for funding my plane fare to Alaska in 2018 and 2019, thus facilitating my visits to the Upper Tanana area.

\textit{Baasee'} to David Engles for discussing these topics with me during our trip to the airport and to Siri Tuttle for providing comments on an earlier draft of this paper. \textit{Marci cho} to Dagmar Jung for the Dëne Sųłiné example. \textit{Merci} to the participants in the workshop `Lest we miss them' at the at the VIII Syntax of the World's Languages conference, September 3–5, 2018, in Paris, France for their comments and questions. \textit{Thank you} to two anonymous reviewers and to the editors, whose comments and questions benefitted the paper considerably.

While I could not have written this paper without the individuals named above, all errors of fact, representation, or interpretation are my own.

\printbibliography[heading=subbibliography,notkeyword=this]
\end{document}
